\chapter{The parameter-derivative tower}
\label{ch:08-differentiation}

Differentiating in the Hurwitz parameter moves the spectral argument upward.
The resulting coefficients are elementary symmetric functions of reciprocal
integers, so harmonic numbers control every rung. The functional equation
then pairs this tower with the integrated tower of the preceding chapter.
Explicit character tables show how experiments can be replaced by uniform
derivations.

\section{The uniform master identity}\label{tower:sec:master}

\begin{theorem}[Harmonic-number master identity]\label{tower:thm:master}
For all $k\ge1$ and $n\ge0$,
\begin{equation}\label{tower:eq:masterN}
  \gamma_n^{(k)}(a)=(-1)^{n+k}\,k!\,n!\sum_{i+j=n} e_i(k)\,
   \frac{\zeta^{(j)}(k{+}1,a)}{j!},
\end{equation}
where $e_i(k)$ is the $i$-th elementary symmetric function of
$\{1,\tfrac12,\dots,\tfrac1k\}$; in particular $e_0=1$, $e_1=\h{k}$,
$e_2=\tfrac12(\h{k}^2-\h{k}^{(2)})$. The first two rows are
\begin{align}
  \gamma_1^{(k)}(a)&=(-1)^{k+1}k!\,\bigl[\zeta'(k{+}1,a)+\h{k}\,\zeta(k{+}1,a)\bigr],
   \label{tower:eq:master1}\\
  \gamma_2^{(k)}(a)&=(-1)^{k}k!\,\bigl[\zeta''(k{+}1,a)+2\h{k}\,\zeta'(k{+}1,a)
   +(\h{k}^2-\h{k}^{(2)})\,\zeta(k{+}1,a)\bigr].\label{tower:eq:master2}
\end{align}
\end{theorem}

\begin{proof}
From $\partial_a\zeta(s,a)=-s\,\zeta(s{+}1,a)$, iteration gives
$\partial_a^k\zeta(s,a)=(-1)^k (s)_k\,\zeta(s{+}k,a)$ with the rising factorial
$(s)_k=s(s{+}1)\cdots(s{+}k{-}1)$. For $k\ge1$ this annihilates the pole term in
\eqref{stieltjes:eq:laurent}, so
\[
  \sum_{n\ge0}\frac{(-1)^n}{n!}\,\gamma_n^{(k)}(a)\,(s-1)^n
   =(-1)^k(s)_k\,\zeta(s{+}k,a).
\]
Set $s=1+\varepsilon$. Then
$(s)_k=(1+\varepsilon)_k=\prod_{m=1}^k(m+\varepsilon)=k!\prod_{m=1}^k(1+\varepsilon/m)
 =k!\sum_{i\ge0}e_i(k)\,\varepsilon^i$,
the generating identity for the elementary symmetric functions of $\{1/m\}$,
while $\zeta(s{+}k,a)=\zeta(k{+}1{+}\varepsilon,a)=\sum_{j\ge0}\zeta^{(j)}(k{+}1,a)\,\varepsilon^j/j!$
is analytic at $\varepsilon=0$ because $k{+}1\ge2$. Comparing coefficients of
$\varepsilon^n$ yields \eqref{tower:eq:masterN}; $n=1,2$ give
\eqref{tower:eq:master1}--\eqref{tower:eq:master2}.
\end{proof}

Theorem~\ref{tower:thm:master} is the harmonic-number repackaging of Coffey's
all-orders Stirling-number ladder~\cite{stieltjes:Coffey}: the change of basis from
the unsigned Stirling numbers $\bigl|s(k{+}1,i{+}1)\bigr|$ to the symmetric
functions $e_i(k)$ turns the coefficients into $\h{k}$, $\h{k}^2-\h{k}^{(2)}$,
\dots. The specializations $\gamma_1'=\zeta'(2,a)+\zeta(2,a)$,
$\gamma_1''=-2\zeta'(3,a)-3\zeta(3,a)$,
$\gamma_2'=-[\zeta''(2,a)+2\zeta'(2,a)]$ and
$\gamma_2''=2\zeta''(3,a)+6\zeta'(3,a)+2\zeta(3,a)$ are the $s=2,3$ layers used
below. Identity~\eqref{tower:eq:masterN} was verified against numerical $k$-th
parameter-derivatives of $\gamma_1,\gamma_2$ for $k=1,\dots,5$.

\section{Differentiation and the contact-term correction}
\label{stconv:sec:contact}
For $k\ge1$, the ordinary function $\gamma_n^{(k)}(x)$ has a power-log
singularity of order $k+1$. Its canonical Hadamard finite part on $\stconvT$
is defined by
\begin{equation}\label{stconv:eq:Fdefinition}
 \langle F_{n,k},\phi\rangle
 =\stconvCT_{\epsilon\downarrow0}
       \int_\epsilon^1\gamma_n^{(k)}(x)\phi(x)\dd x.
\end{equation}
Here $\stconvCT$ means the coefficient independent of $\epsilon$ after the
finite expansion in negative powers of $\epsilon$ and nonconstant powers
of $\log\epsilon$ is removed. The remaining error tends to zero.
This convention uses the coordinate $x\in(0,1)$ and no finite counterterm.
It agrees with Definition~\ref{stconv:def:G} when $k=0$.

\begin{lemma}[Higher Hurwitz poles as distributions]\label{stconv:lem:higherpoles}
For each $k\ge1$, as $z\to0$,
\begin{equation}\label{stconv:eq:higherpole}
 Z(1+k-z)=\frac{(-1)^k}{k!}\frac{\delta_0^{(k)}}z
 +\sum_{j\ge0}\frac{(-1)^jz^j}{j!}
           \stconvFp_{\stconvT}[\zeta^{[j]}(k+1,x)].
\end{equation}
The finite parts in this formula have the cutoff convention
\eqref{stconv:eq:Fdefinition}.
\end{lemma}
\begin{proof}
Split $\zeta(1+k-z,x)=x^{-k-1+z}+\zeta(1+k-z,1+x)$.
Subtract the Taylor polynomial of $\phi$ at zero through degree $k$.
Its remainder gives an analytic integral near $z=0$. The subtracted
monomials give terms
$\phi^{(r)}(0)/(r!(z+r-k))$, $0\le r\le k$; only $r=k$ has a pole
at zero. Since $\langle\delta_0^{(k)},\phi\rangle=(-1)^k\phi^{(k)}(0)$,
its residue is the one stated. Expanding each regular term and each
analytic remainder is exactly the constant-term prescription for the
cutoff integrals. This proves the entire Laurent expansion.
\end{proof}

\begin{theorem}[All-index differentiation anomaly]\label{stconv:thm:anomaly}
For $n\ge0$ and $k\ge1$,
\begin{equation}\label{stconv:eq:anomaly}
 D^kG_n=F_{n,k}+(-1)^{n+1}n!e_{n+1}(k)\delta_0^{(k)}.
\end{equation}
In particular, there is no anomaly when $n\ge k$.
\end{theorem}
\begin{proof}
The meromorphic distribution family satisfies
\begin{equation}\label{stconv:eq:DZ}
 D^kZ(s)=(-1)^k(s)_kZ(s+k).
\end{equation}
This follows directly from its Fourier coefficients and the Gamma
recurrence; it also follows by repeated integration by parts in a
left half-plane and meromorphic continuation. Put $s=1-z$ and expand
\[
 (1-z)_k=k!\prod_{r=1}^k(1-z/r)
        =k!\sum_{i=0}^k(-1)^ie_i(k)z^i.
\]
On the left of~\eqref{stconv:eq:DZ}, the pole is $\delta_0^{(k)}/z$ and the
regular coefficient of $z^n$ is $D^kG_n/n!$. On the right, multiplying
the regular part of Lemma~\ref{stconv:lem:higherpoles} by the polynomial gives
$F_{n,k}/n!$: this is the finite part of the ordinary pointwise
parameter-derivative identity. Multiplying its polar part by that same
polynomial gives, in addition, the regular coefficient
$(-1)^{n+1}e_{n+1}(k)\delta_0^{(k)}$. Comparison proves the formula.
\end{proof}

The first examples are
\begin{align}
 DG_0&=F_{0,1}-\delta_0',\label{stconv:eq:anomalylow}\\
 D^2G_0&=F_{0,2}-\tfrac32\delta_0'',\notag\\
 D^2G_1&=F_{1,2}+\tfrac12\delta_0''.\notag
\end{align}
The point-supported terms disappear upon restriction to $(0,1)$, so
ordinary differentiation there cannot detect them. But convolving
$\delta_0^{(k)}$ with another distribution differentiates that other
factor everywhere. Therefore these terms cannot be discarded before
convolution, even when the desired final identity is off the singular
point.

The pointwise harmonic master identity \eqref{tower:eq:masterN} remains
valid. The new contact term governs its canonical periodic extension:
point-supported derivatives vanish on the open interval but contribute
after convolution.

\subsection{Convolutions of arbitrary differentiated Stieltjes functions}
The same argument gives a closed law for all derivative indices, not
only the zeroth Stieltjes function. Put $F_{n,0}=G_n$ and define
\[
 a_{n,k}=(-1)^{n+1}n!e_{n+1}(k)\quad(k\ge1),\qquad a_{n,0}=0.
\]
Thus $F_{n,k}=D^kG_n-a_{n,k}\delta_0^{(k)}$ when $k\ge1$.

\begin{corollary}[The full differentiated convolution calculus]
\label{stconv:cor:fulltower}
Suppose $n,m,k,l\ge0$, $r=k+l\ge1$, and the exact rule from
Theorem~\ref{stconv:thm:product} is
\[
 G_n*G_m=c_{-1}\stconvId+\sum_{j=0}^{n+m+1}c_jG_j.
\]
Then
\begin{align}
 F_{n,k}*F_{m,l}={}&\sum_{j=0}^{n+m+1}c_jD^rG_j
             -a_{n,k}D^rG_m-a_{m,l}D^rG_n\label{stconv:eq:fulltower}\\
 &+(c_{-1}+a_{n,k}a_{m,l})\delta_0^{(r)}.\notag
\end{align}
Away from the integer point this reduces to
\[
 \sum_{j=0}^{n+m+1}c_j\gamma_j^{(r)}(x)
          -a_{n,k}\gamma_m^{(r)}(x)-a_{m,l}\gamma_n^{(r)}(x).
\]
In particular, every such convolution has a finite reduction to
$\zeta^{[j]}(r+1,x)$ with $0\le j\le n+m+1$.
\end{corollary}
\begin{proof}
Multiply the two identities $F=D G-a\delta$ at their indicated
orders, interpreting the zero-order case through $a_{n,0}=0$.
The product of the first terms is $D^r(G_n*G_m)$, and
$D^r\stconvId=\delta_0^{(r)}$. Restricting off zero and applying
\eqref{tower:eq:masterN} proves the last assertion.
\end{proof}

As an explicit higher-index example, the contact term vanishes for
$F_{1,1}=DG_1$. Hence
\begin{align}
 (F_{1,1}*F_{1,1})(x)
 ={}&-2\zeta'''(3,x)-9\zeta''(3,x)\label{stconv:eq:gamma1primepair}\\
 &+[4\zeta(2)-6]\zeta'(3,x)
   +[6\zeta(2)+4\zeta(3)]\zeta(3,x).\notag
\end{align}
This follows by twice differentiating~\eqref{stconv:eq:product11} and using
\eqref{tower:eq:masterN}. It illustrates that the full tower still
closes at one spectral integer even when higher Stieltjes indices occur.

\section{A two-jet collapse for every polygamma convolution}
\label{stconv:sec:poly}
Let $H_0=0$ and $H_k=\sum_{j=1}^k1/j$ for $k\ge1$. Define
\[
 \Psi_k=\stconvFp_{\stconvT}\psi^{(k)}\qquad(k\ge0)
\]
using the same canonical cutoff convention. Since $\gamma_0=-\psi$,
Theorem~\ref{stconv:thm:anomaly} gives
\begin{equation}\label{stconv:eq:PsiD}
 \Psi_0=-G_0,\qquad
 \Psi_k=-D^kG_0-H_k\delta_0^{(k)}\quad(k\ge1).
\end{equation}

\begin{theorem}[Polygamma convolution reduction]\label{stconv:thm:polygamma}
Let $k,l\ge0$ and $r=k+l\ge1$. Then
\begin{align}
 \Psi_k*\Psi_l={}&2D^rG_1+(H_k+H_l)D^rG_0\label{stconv:eq:polyfull}\\
 &+[H_kH_l-\zeta(2)]\delta_0^{(r)}.\notag
\end{align}
Consequently, for $0<x<1$,
\begin{equation}\label{stconv:eq:polyscalar}
 (\Psi_k*\Psi_l)(x)=(-1)^{r+1}r!
 \left[2\zeta'(r+1,x)+(2H_r-H_k-H_l)\zeta(r+1,x)\right].
\end{equation}
For $k=l=0$, the appropriate separate formula is
$\Psi_0*\Psi_0=2G_1-\zeta(2)\stconvId$.
\end{theorem}
\begin{proof}
Convolving~\eqref{stconv:eq:PsiD} and using that derivatives may be placed on
either factor gives
\[
 \Psi_k*\Psi_l=D^r(G_0*G_0)
 +(H_k+H_l)D^rG_0+H_kH_l\delta_0^{(r)}.
\]
The same expression is valid when one index is zero, since its harmonic
number vanishes. Apply~\eqref{stconv:eq:product00} and
$D^r\stconvId=\delta_0^{(r)}$. Restriction to $0<x<1$ gives
$2\gamma_1^{(r)}+(H_k+H_l)\gamma_0^{(r)}$.
Now use~\eqref{tower:eq:masterN} with $n=0,1$:
\[
 \gamma_1^{(r)}=(-1)^{r+1}r![\zeta'(r+1,x)+H_r\zeta(r+1,x)],
 \quad \gamma_0^{(r)}=(-1)^rr!\zeta(r+1,x).
\]
This proves~\eqref{stconv:eq:polyscalar}.
\end{proof}

For example,
\begin{align}
 (\Psi_0*\Psi_1)(x)&=2\zeta'(2,x)+\zeta(2,x),\label{stconv:eq:polyexamples}\\
 (\Psi_1*\Psi_1)(x)&=-4\zeta'(3,x)-2\zeta(3,x),\notag\\
 (\Psi_0*\Psi_2)(x)&=-4\zeta'(3,x)-3\zeta(3,x).\notag
\end{align}
Although the last two convolutions have the same total derivative order,
they have different coefficients. Confusing $\Psi_k$ with the $k$th
derivative of $\Psi_0$ would incorrectly erase this distinction.

\subsection{A convergent trigamma identity with every subtraction visible}
For $0<x<1$, set $h=x/2$ and $f(t)=\psi'(t)$. Then
\begin{align}
 \mathcal T(x)={}&2\int_0^h\left[
       \frac{f(x-t)-f(x)+tf'(x)}{t^2}
        +\psi'(1+t)f(x-t)\right]\dd t\label{stconv:eq:trigamma-integral}\\
 &+\int_x^1f(t)f(1+x-t)\dd t
      -\frac{2f(x)}h-2f'(x)\log h\notag
\end{align}
satisfies
\begin{equation}\label{stconv:eq:trigamma-value}
 \mathcal T(x)=-4\zeta'(3,x)-2\zeta(3,x).
\end{equation}
Indeed, the shift relation $\psi'(t)=t^{-2}+\psi'(1+t)$ and symmetry
about $t=x/2$ reduce the first singular convolution integral to twice
its left half. Subtracting the constant and linear Taylor terms of
$f(x-t)$ leaves a bounded quotient. The constant terms of
$\int_\epsilon^ht^{-2}\dd t$ and
$\int_\epsilon^ht^{-1}\dd t$ are $-1/h$ and $\log h$,
respectively. This proves that~\eqref{stconv:eq:trigamma-integral} is exactly
the canonical finite-part convolution, so Theorem~\ref{stconv:thm:polygamma}
applies. Every integral in~\eqref{stconv:eq:trigamma-integral} is ordinary and
absolutely convergent.



\section{Contact terms in translated derivative products}
\label{stconv:sec:translated-contact}
For $r\ge1$ let $h_{r,m}=(-1)^m m!e_{m+1}(r)$, and put $h_{0,m}=0$.
Thus \eqref{stconv:eq:anomaly} reads
$F_{m,r}=D^rG_m+h_{r,m}\delta_0^{(r)}$.
\begin{theorem}[Translated contact law]\label{shj:contactthm}
For $r,k,m,n\ge0$, $N=r+k$, and $b=1-a$ with $0<a<1$,
\begin{align}\label{shj:generalcontact}
 &\FP\int_0^1\gamma_m^{(r)}(x)\gamma_n^{(k)}(\{x+a\})\,dx\notag\\
 &\quad=(-1)^r J_{mn}^{(N)}(a)
 +(-1)^k h_{k,n}\gamma_m^{(N)}(b)
 +(-1)^r h_{r,m}\gamma_n^{(N)}(a).
\end{align}
\end{theorem}
\begin{proof}
The local finite parts use the constant-term convention
\eqref{stconv:eq:Fdefinition}. At nonzero shift the singular supports
of the two factors are disjoint, so a partition into disjoint endpoint
neighborhoods defines the product unambiguously. Substitute the two contact
identities. Periodic integration by parts gives the first term. The
shifted delta derivative acts at $x=b$, giving
$(-1)^k h_{k,n}\gamma_m^{(N)}(b)$, and the unshifted one acts at zero,
giving $(-1)^r h_{r,m}\gamma_n^{(N)}(a)$. The product of the two
point-supported terms vanishes because their supports are distinct.
This proves the formula with the exact scale-one endpoint normalization.
\end{proof}
\subsection{Every polygamma product in first Hurwitz derivatives}
For $r,k\ge0$ define
\[
\Pi_{r,k}(a)=\FP\int_0^1\psi^{(r)}(x)\psi^{(k)}(\{x+a\})\dd x.
\]
\begin{corollary}[Polygamma closure]\label{shj:polyclosure}
For $N=r+k\ge1$,
\begin{align}
\Pi_{r,k}(a)={}&(-1)^r\big[\gamma_1^{(N)}(a)-H_r\psi^{(N)}(a)\big]\notag\\
&+(-1)^k\big[\gamma_1^{(N)}(b)-H_k\psi^{(N)}(b)\big].
\label{shj:polyclosuregamma}
\end{align}
Equivalently,
\begin{align}
\Pi_{r,k}(a)=N!\bigg\{&(-1)^{k+1}\big[\zeta^{[1]}(N+1,a)
+(H_N-H_r)\zeta(N+1,a)\big]\notag\\
&+(-1)^{r+1}\big[\zeta^{[1]}(N+1,b)
+(H_N-H_k)\zeta(N+1,b)\big]\bigg\}.
\label{shj:polyclosurezeta}
\end{align}
The case $N=0$ is \eqref{shj:psipsi}.
\end{corollary}
\begin{proof}
Take $m=n=0$ in \eqref{shj:generalcontact}, use $\gamma_0^{(r)}=-\psi^{(r)}$, $h_{r,0}=H_r$, and differentiate \eqref{shj:J00}. This proves \eqref{shj:polyclosuregamma}. The pointwise identities
\begin{align*}
\psi^{(N)}(x)&=(-1)^{N+1}N!\zeta(N+1,x),\\
\gamma_1^{(N)}(x)&=(-1)^{N+1}N!
\big[\zeta^{[1]}(N+1,x)+H_N\zeta(N+1,x)\big]
\end{align*}
then give \eqref{shj:polyclosurezeta}. The second follows directly by expanding
$\partial_x^N\zeta(1+u,x)=(-1)^N(1+u)_N\zeta(N+1+u,x)$ at $u=0$; these pointwise ladders are also part of the repository's existing development \cite{stconv:shifted-report}.
\end{proof}

Two instructive special cases are
\begin{align}
\Pi_{0,1}(a)&=\zeta^{[1]}(2,a)+\zeta(2,a)-\zeta^{[1]}(2,b)
=J_{00}'(a)+\psi'(b),\label{shj:01}\\
\Pi_{1,1}(a)&=2\zeta^{[1]}(3,a)+2\zeta^{[1]}(3,b)
+\zeta(3,a)+\zeta(3,b).\label{shj:11}
\end{align}
The term $+\psi'(b)$ in \eqref{shj:01} is an explicit counterexample to the unqualified rule ``differentiate the finite-part identity under the integral.'' It is not a counterexample to ordinary differentiation under a genuinely dominated integral.

For reproducibility, let $I_{r,k}(\epsilon)$ be the sum of the two raw cutoff integrals, with arguments split as in \eqref{shj:Jcutoff}. Then
\begin{align}
\Pi_{0,1}(a)=\lim_{\epsilon\downarrow0}\left[
I_{0,1}(\epsilon)-\frac{\psi(b)}\epsilon
-\big(\psi'(a)-\psi'(b)\big)\log\epsilon\right],\label{shj:cut01}\\
\Pi_{1,1}(a)=\lim_{\epsilon\downarrow0}\left[
I_{1,1}(\epsilon)-\frac{\psi'(a)+\psi'(b)}\epsilon
+\big(\psi''(a)+\psi''(b)\big)\log\epsilon\right].\label{shj:cut11}
\end{align}
These counterterms can be derived directly from $\psi(x)=-1/x-\gamma+O(x)$ and $\psi'(x)=x^{-2}+O(1)$. They give numerical tests of the contact law that do not implement distributional differentiation.


The full default replays are preserved under \src{verification/eighth-replay/}.
The direct cutoff polygamma calculations check the contact correction
without implementing distributional differentiation. The all-index law
rests on the Laurent residues and the proof above. The ordinary pointwise
master identity has not been retracted or altered.

\section{Dilation, distribution traces and exact logarithmic contacts}
\label{dilfp:sec:main}
The contact theorem answers an extension question suggested by the
correlation reports: classical distribution rows need a normalization
correction when interpreted as finite-part identities. Fourier traces
make the correction computable at every Stieltjes and derivative order.
This also describes how dilation acts on the one-generator convolution
algebra, without any arithmetic-independence assumption.

For a positive integer $q$, define the trace $\mathcal S_q$ and covering
pullback $\mathcal P_q$ on periodic distributions by
\[
 \widehat{\mathcal S_qT}(k)=q\widehat T(qk),\qquad
 \widehat{\mathcal P_qT}(k)=
 \begin{cases}\widehat T(k/q),&q\mid k,\\0,&q\nmid k.\end{cases}
\]
For ordinary periodic functions these are
$\mathcal S_qf(x)=\sum_{j=0}^{q-1}f((x+j)/q)$ and
$\mathcal P_qf(x)=f(\{qx\})$. Put $\ell_q=\log q$ and
$\mathcal L_q=q^{-1}\mathcal S_q$.

\begin{theorem}[The trace translates the convolution generator]
\label{dilfp:thm:trace}
For every $m\ge0$,
\begin{equation}\label{dilfp:eq:trace}
 \mathcal S_qG_m
 =q\sum_{j=0}^m\binom mj(-\ell_q)^{m-j}G_j
   +\frac{q(-\ell_q)^{m+1}}{m+1}\stconvId.
\end{equation}
The normalized trace $\mathcal L_q$ is a unital homomorphism of the
mean-zero convolution algebra. On its generator $X=G_0$ it acts by
$X\mapsto X-\ell_q\stconvId$. These actions compose as
$\mathcal L_q\mathcal L_r=\mathcal L_{qr}$.
\end{theorem}
\begin{proof}
The Fourier definition gives
$\mathcal S_qR_\alpha=q^{1-\alpha}R_\alpha$. Apply this to the
Laurent family \eqref{stconv:eq:ZLaurent}:
\[
 \mathcal S_q\left(\frac{\stconvId}{z}
                  +\sum_{m\ge0}G_m\frac{z^m}{m!}\right)
 =q e^{-\ell_qz}\left(\frac{\stconvId}{z}
                  +\sum_{m\ge0}G_m\frac{z^m}{m!}\right).
\]
The regular coefficients give \eqref{dilfp:eq:trace}, including the
term produced by the pole. On Fourier coefficients,
$\mathcal L_q(T*S)=(\mathcal L_qT)*(\mathcal L_qS)$ and
$\mathcal L_q\stconvId=\stconvId$. The generator formula is the
$m=0$ case, and repeated traces evaluate a Fourier coefficient at $qrk$.
\end{proof}

\begin{theorem}[Finite part does not commute with covering pullback]
\label{dilfp:thm:cover}
Let $\operatorname{Fp}_x\gamma_m(\{qx\})$ use the constant term in
the unscaled local coordinate $x-j/q$ at each of the $q$ singularities,
with right-hand cutoffs tending independently to zero. Then
\begin{equation}\label{dilfp:eq:cover}
 \mathcal P_qG_m
 =\operatorname{Fp}_x\gamma_m(\{qx\})
 +\frac{\ell_q^{m+1}}{q(m+1)}
                    \sum_{j=0}^{q-1}\delta_{j/q}.
\end{equation}
In particular the raw-coordinate finite part has mean
$-\ell_q^{m+1}/(m+1)$, while the covering pullback has mean zero.
\end{theorem}
\begin{proof}
On the $j$th segment substitute $y=q(x-j/q)$. Pulling back the
cutoff prescription \eqref{stconv:eq:Gcutoff} gives its local
counterterm
\[
 \frac{\log^{m+1}(q\epsilon)}{q(m+1)}\phi(j/q).
\]
The raw-coordinate constant-term prescription keeps only the divergent
powers of $\log\epsilon$ in this expression. Their difference is its
constant coefficient, $\ell_q^{m+1}\phi(j/q)/(q(m+1))$.
All Taylor remainder integrals converge as in the definition of $G_m$.
Summing the segments proves \eqref{dilfp:eq:cover}. The Fourier
pullback preserves the zeroth coefficient, and pairing the contact sum
with the constant function proves the mean statement.
\end{proof}

For the trace itself, put $c_{m,q}=-q(-\ell_q)^{m+1}/(m+1)$.
The classical pointwise multiplication row and its raw-coordinate
extension read
\[
 \operatorname{Fp}_x[\mathcal S_q\gamma_m]
 =q\sum_{j=0}^m\binom mj(-\ell_q)^{m-j}G_j+c_{m,q}\,1.
\]
Comparing with \eqref{dilfp:eq:trace} gives the separate identity
\begin{equation}\label{dilfp:eq:trace-commutator}
 \mathcal S_qG_m
 =\operatorname{Fp}_x[\mathcal S_q\gamma_m]-c_{m,q}\delta_0.
\end{equation}
The logarithmic term thus has different signs and support for covering
pullback and trace; substituting one convention for the other is invalid.

\begin{theorem}[Every differentiated trace and its contact]
\label{dilfp:thm:derivative}
Use $F_{m,r}$ and $a_{m,r}=(-1)^{m+1}m!e_{m+1}(r)$ from the
contact theorem, and let $r\ge1$. Then
\begin{align}\label{dilfp:eq:derivative}
 \mathcal S_qF_{m,r}
 ={}&q^{r+1}\sum_{j=0}^m\binom mj(-\ell_q)^{m-j}F_{j,r}
                         +q^{r+1}\kappa_{m,r}(\ell_q)\delta_0^{(r)},\notag\\
 \kappa_{m,r}(\ell)
 ={}&\frac{(-\ell)^{m+1}}{m+1}
       +\sum_{j=0}^m\binom mj(-\ell)^{m-j}a_{j,r}-a_{m,r}.
\end{align}
The first sum is exactly
$\operatorname{Fp}_x[\mathcal S_q\gamma_m^{(r)}]$.
For the canonical polygamma finite parts, the particularly simple
all-order specialization is
\begin{equation}\label{dilfp:eq:polygamma}
 \mathcal S_q\Psi_r=q^{r+1}
                   (\Psi_r+\ell_q\delta_0^{(r)}),\qquad r\ge1.
\end{equation}
\end{theorem}
\begin{proof}
Fourier coefficients give
$\mathcal S_qD^r=q^rD^r\mathcal S_q$ and
$\mathcal S_q\delta_0^{(r)}=q^{r+1}\delta_0^{(r)}$.
Apply these identities to $F_{m,r}=D^rG_m-a_{m,r}\delta_0^{(r)}$,
use \eqref{dilfp:eq:trace}, and then replace each
$D^rG_j$ by $F_{j,r}+a_{j,r}\delta_0^{(r)}$.
Because $D^r\stconvId=\delta_0^{(r)}$, the contact coefficient is
precisely the displayed $\kappa_{m,r}$. Differentiating the classical
pointwise multiplication row $r$ times proves the raw-coordinate
interpretation of the first sum. Finally $\Psi_r=-F_{0,r}$ and
$\kappa_{0,r}=-\ell_q$, proving \eqref{dilfp:eq:polygamma}.
\end{proof}

These identities answer the normalization question for integer covering
maps and their distribution traces at every index. They also show that
the polynomial convolution algebra is preserved by the normalized trace.
Twisted character projections and products at several independent shifted
singularities still require separate calculations. Exact Fourier-polynomial
and local logarithmic-counterterm controls, with a corruption check, are
recorded in \src{verification/certify_stieltjes_dilation.py}.


\section{The uniform rational-argument landscape}\label{tower:sec:landscape}

Three structural facts hold uniformly in $k$; all follow from
$\sum_{p}\chi(p)\,\zeta(s,p/q)=q^{s}L(s,\chi)$ (and its multiplicative cousin)
differentiated in $s$, fed through Theorem~\ref{tower:thm:master}.

\begin{proposition}[Character coordinate]\label{tower:prop:char}
For a primitive nonprincipal Dirichlet character $\chi$ modulo $q>1$ and all $k\ge1$,
\begin{equation}\label{tower:eq:char}
  \sum_{p=1}^{q-1}\chi(p)\,\gamma_1^{(k)}(p/q)
   =(-1)^{k+1}k!\,q^{k+1}\bigl[L'(k{+}1,\chi)+(\h{k}+\ln q)\,L(k{+}1,\chi)\bigr].
\end{equation}
\end{proposition}

\begin{proposition}[Distribution row]\label{tower:prop:dist}
For $d\ge2$ and all $k\ge1$, with $\psi^{(k)}(a)=(-1)^{k+1}k!\,\zeta(k{+}1,a)$,
\begin{equation}\label{tower:eq:dist}
  \sum_{j=0}^{d-1}\gamma_1^{(k)}\!\Bigl(\tfrac{a+j}{d}\Bigr)
   =d^{\,k+1}\gamma_1^{(k)}(a)+d^{\,k+1}\ln d\;\psi^{(k)}(a).
\end{equation}
\end{proposition}

\begin{proof}[Proof of \eqref{tower:eq:char} and \eqref{tower:eq:dist}]
Apply the character sum to \eqref{tower:eq:master1} for
\eqref{tower:eq:char}, and the divisor sum for \eqref{tower:eq:dist}.
For the character identity, use
$\sum_p\chi(p)\zeta(s,p/q)=q^sL(s,\chi)$ and its $s$-derivative
$q^s\ln q\,L(s,\chi)+q^sL'(s,\chi)$ at $s=k{+}1$; collect. For \eqref{tower:eq:dist},
use $\sum_{j<d}\zeta(s,(a{+}j)/d)=d^s\zeta(s,a)$ and its $s$-derivative, and
recognize $(-1)^{k+1}k!\,\zeta(k{+}1,a)=\psi^{(k)}(a)$.
\end{proof}

\begin{corollary}[Uniform formal distribution ranks]\label{tower:thm:rank}
For every $q\ge2$ and every parameter derivative order $k\ge1$, the
specified distribution presentation, with the endpoint and lower-index
data treated as background, has residual dimension $\varphi(q)-1$.
This is a formal quotient statement, not arithmetic independence of the
evaluated constants.
\end{corollary}
The original finite rank experiments suggested this law. The following
polynomial character normal form proves it without extrapolation, also
at arbitrary weights and in finite jet algebras.
\section{Weighted distribution modules and exact rational-grid ranks}
\label{dist:sec:distribution}

The distribution relations in the Stieltjes drafts have a simple
representation-theoretic structure. The relevant dimension is the
dimension of a formal quotient by specified relations. The classical
theory of universal distributions provides the broader setting
\cite{dist:Kubert,dist:Ouyang}. Here we give an explicit character normal form,
including its behavior at exceptional weights and in finite Taylor
jets. This proves the two all-denominator rank statements suggested by
the source computations, without making an independence assertion about
evaluated special values.

\subsection{The relation module and its complete character decomposition}

Fix \(q\geq2\), and put
\[
 G_q=\frac1q\Z/\Z,\qquad U(q)=(\Z/q\Z)^\times.
\]
Let \(\mathcal A\) be a commutative unital \(\C\)-algebra and choose
arbitrary \(t_p\in\mathcal A\), one for each prime \(p\mid q\).
The choices may vanish or be nilpotent. Let
\[
 \mathcal V_q=\bigoplus_{x\in G_q}\mathcal A[x].
\]
For \(p\mid q\) and \(x\in pG_q\), define the prime distribution row
\begin{equation}
 D_p(x)=\sum_{\substack{y\in G_q\\py=x}}[y]-t_p[x].
 \label{dist:eq:dist-prime-row}
\end{equation}
Write \(\mathcal R_q\) for the \(\mathcal A\)-span of these rows and
\(\mathcal U_q=\mathcal V_q/\mathcal R_q\).
The unit group acts by \(u[x]=[ux]\), and each relation map is
equivariant.

For \(d\mid q\), set \(t_d=\prod_{p\mid d}t_p^{v_p(d)}\).
Rows for composite divisors introduce no additional relations. Indeed,
if \(ab\mid q\) and \(x\in abG_q\), then
\begin{equation}
 D_{ab}(x)=
 \sum_{\substack{z\in G_q\\az=x}}D_b(z)+t_bD_a(x).
 \label{dist:eq:dist-composite-row}
\end{equation}
Every \(z\) in this sum belongs to \(bG_q\), so the rows on the right
are legitimate grid rows. Expanding the two sums proves the identity.
Induction on the number of prime factors of \(d\), counted with
multiplicity, proves the assertion for every divisor.

Let \(\chi\) run through all characters of \(U(q)\), and let \(f=f_\chi\)
be the conductor of its primitive inducing character \(\chi^*\).
For \(f\mid d\mid q\), let \(\chi_d\) denote the induced character
modulo \(d\), and define
\begin{equation}
 Y_{d,\chi}=\sum_{a\in U(d)}\chi_d(a)[a/d].
 \label{dist:eq:dist-character-vector}
\end{equation}
For the principal character at \(d=1\), this means
\(Y_{1,\mathbf1}=[0]\). These are formal vectors; no numerical
\(L\)-value is being declared independent.

\begin{lemma}[All character components of the rows]
\label{dist:lem:dist-all-components}
The \(\chi^{-1}\)-eigenspace of \(\mathcal V_q\) is free over
\(\mathcal A\), with basis
\[
 \{Y_{d,\chi}: f_\chi\mid d\mid q\}.
\]
Its relation submodule is generated exactly by the edges
\begin{equation}
 \begin{aligned}
 Y_{pd,\chi}-t_pY_{d,\chi},
     &\qquad p\mid d,\\
 Y_{pd,\chi}-(t_p-\chi^*(p))Y_{d,\chi},
     &\qquad p\nmid d,
 \end{aligned}
 \qquad f_\chi\mid d,\quad pd\mid q.
 \label{dist:eq:dist-character-edges}
\end{equation}
In particular, there are no additional character relations arising
from denominators smaller than, or not divisible by, \(f_\chi\).
\end{lemma}

\begin{proof}
The elements of \(G_q\) of exact denominator \(d\) form a single
\(U(q)\)-orbit, naturally identified with \(U(d)\). Reduction
\(U(q)\to U(d)\) is surjective. The character \(\chi\) factors through
this map exactly when \(f_\chi\mid d\); otherwise its projection on
that orbit is zero. More explicitly, the idempotent
\[
 E_\chi=\frac1{\varphi(q)}\sum_{u\in U(q)}\chi(u)u
\]
satisfies, for \(a\in U(d)\),
\begin{equation}
 E_\chi[a/d]=
 \begin{cases}
 \displaystyle\frac{\chi_d(a)^{-1}}{\varphi(d)}Y_{d,\chi},
       &f_\chi\mid d,\\[6pt]
 0,&f_\chi\nmid d.
 \end{cases}
 \label{dist:eq:dist-projector}
\end{equation}
This follows by grouping \(u\) according to its residue modulo \(d\)
and using character orthogonality on the kernel of reduction.
The idempotents are mutually orthogonal and sum to the identity.
All their scalar denominators are units in \(\mathcal A\), so the
argument applies equally to a nonreduced algebra.

For a fixed \(p\mid q\), consider the equivariant map
\[
 D_p:\mathcal A[pG_q]\longrightarrow\mathcal V_q,\qquad
 [x]\longmapsto D_p(x).
\]
The exact-denominator orbits in its domain have \(pd\mid q\).
By \eqref{dist:eq:dist-projector}, their only nonzero \(\chi\)-components
have \(f_\chi\mid d\), and each such component has the single
generator \(Y_{d,\chi}\). Consequently, computing
\(D_p(Y_{d,\chi})\) computes \emph{every} projected relation for
that \(p\); an orbit on which the projector vanishes cannot create
a hidden relation in the target.

If \(p\mid d\), all \(p\) preimages of each \(a/d\), \(a\in U(d)\),
have exact denominator \(pd\). Summing with weights \(\chi_d(a)\)
therefore gives \(Y_{pd,\chi}-t_pY_{d,\chi}\).
If \(p\nmid d\), one preimage has denominator \(d\), while the other
\(p-1\) have denominator \(pd\). In the former preimage,
\(pc\equiv a\pmod d\), so its character coefficient is
\(\chi_d(a)=\chi^*(p)\chi_d(c)\). Thus the sum is
\[
 Y_{pd,\chi}+\chi^*(p)Y_{d,\chi}-t_pY_{d,\chi}.
 \]
Here \(p\nmid f_\chi\), as \(f_\chi\mid d\), so \(\chi^*(p)\)
is a nonzero root of unity. This proves both the edge formulas
and their completeness.
\end{proof}

\begin{theorem}[Polynomial normal form]
\label{dist:thm:dist-normal-form}
For \(f=f_\chi\mid d\mid q\), define
\begin{equation}
 A_{d,\chi}(\mathbf t)=
 \prod_{p\mid f}t_p^{\,v_p(d)-v_p(f)}
 \prod_{\substack{p\mid d\\p\nmid f}}
       t_p^{\,v_p(d)-1}(t_p-\chi^*(p)).
 \label{dist:eq:dist-normal-coefficient}
\end{equation}
Empty products are \(1\), so \(A_{f,\chi}=1\).
Then the full \(\chi\)-relation submodule is
\begin{equation}
 \mathcal R_{q,\chi}
 =\bigoplus_{\substack{f\mid d\mid q\\d\ne f}}
   \mathcal A\bigl(Y_{d,\chi}-A_{d,\chi}Y_{f,\chi}\bigr).
 \label{dist:eq:dist-triangular-relations}
\end{equation}
In particular,
\begin{equation}
 \mathcal U_q\cong
 \bigoplus_{\chi\in\widehat{U(q)}}\mathcal A\,Y_{f_\chi,\chi}
 \cong\mathcal A^{\varphi(q)}.
 \label{dist:eq:dist-free}
\end{equation}
The character normal form of every vector is obtained by replacing
\(Y_{d,\chi}\) with \(A_{d,\chi}Y_{f_\chi,\chi}\).
\end{theorem}

\begin{proof}
The polynomials \(A_{d,\chi}\) satisfy precisely the edge recursions
\eqref{dist:eq:dist-character-edges}: a prime already dividing \(d\)
contributes \(t_p\), and a prime introduced for the first time
outside \(f\) contributes \(t_p-\chi^*(p)\).
The factors for distinct primes commute.

Follow any chain of prime extensions from \(f\) to \(d\).
Successively eliminating the intermediate \(Y\)'s expresses
\(Y_{d,\chi}-A_{d,\chi}Y_{f,\chi}\) as a combination of the
edge relations. Conversely, write an edge as
\(Y_{pd,\chi}-cY_{d,\chi}\), where
\(A_{pd,\chi}=cA_{d,\chi}\). It equals
\[
 \bigl(Y_{pd,\chi}-A_{pd,\chi}Y_{f,\chi}\bigr)
 -c\bigl(Y_{d,\chi}-A_{d,\chi}Y_{f,\chi}\bigr).
 \]
Lemma~\ref{dist:lem:dist-all-components} therefore gives equality of the
two relation spans.

The displayed generators in \eqref{dist:eq:dist-triangular-relations}
are independent: the coefficient of \(Y_{d,\chi}\), for \(d\ne f\),
in a linear combination is exactly the coefficient of its own
generator. Together with \(Y_{f,\chi}\), they form a triangular
basis of the whole character block. Its quotient is consequently
free of rank one. Summing over all \(\varphi(q)\) characters proves
\eqref{dist:eq:dist-free}.
\end{proof}

\begin{remark}
The proof divides by neither \(t_p\) nor \(t_p-\chi^*(p)\).
It proves freeness first over \(\C[t_p:p\mid q]\) and then over
every \(\C\)-algebra obtained by specialization. Exceptional weights
can therefore make a preferred coordinate vanish, but cannot change
the rank of the full quotient. This distinction is essential when
Euler factors are used to solve a relation system.
\end{remark}

\subsection{Anchors, reflection, and the ranks in the drafts}

In this subsection take \(\mathcal A=\C\).
In a Hurwitz-zeta grid, use the representatives \(0<a\leq1\):
the group element \(0\) represents \(a=1\), not the singular
Hurwitz parameter \(a=0\). Treating the endpoint as known means
imposing the homogeneous anchor \([0]=0\). It removes exactly the
principal generator \(Y_{1,\mathbf1}\). Hence
\begin{equation}
 \dim_\C\mathcal V_q/(\mathcal R_q+\C[0])=\varphi(q)-1.
 \label{dist:eq:dist-anchor-dimension}
\end{equation}
There are \(q-1\) remaining grid variables, so their distribution
coefficient matrix has rank \(q-\varphi(q)\). At rational weights
\(t_p\), the matrix has rational entries and this is also its rank
over \(\Q\).

Let \(J[x]=[-x]\). On the \(\chi\)-block, \(J\) acts as
\(\chi(-1)\). Thus the additional homogeneous reflection rows
\begin{equation}
 [x]+(-1)^k[-x]=0,\qquad x\ne0,
 \label{dist:eq:dist-reflection-rows}
\end{equation}
retain exactly the characters with
\(\chi(-1)=(-1)^{k+1}\). After anchoring, the omission of \(x=0\)
in \eqref{dist:eq:dist-reflection-rows} has no effect.

\begin{corollary}[All-denominator formal ranks]
\label{dist:cor:dist-all-ranks}
For \(q\geq3\) and every integer \(k\geq0\), the formal quotient
by distribution, endpoint, and the reflection rows
\eqref{dist:eq:dist-reflection-rows} has dimension
\begin{equation}
 r_{q,k}=
 \begin{cases}
 \varphi(q)/2-1,&k\ \text{odd},\\
 \varphi(q)/2,&k\ \text{even}.
 \end{cases}
 \label{dist:eq:dist-parity-dimension}
\end{equation}
This conclusion holds for arbitrary prime weights.
In particular it proves the all-\(q\), all-\(k\) formal rank law
for the negative Hurwitz-zeta derivative ladder.
The formal distribution-only rank in the Stieltjes parameter
derivative tower is \(q-\varphi(q)\) for every \(q\geq2\),
every parameter derivative order \(k\geq1\), and every Stieltjes
index \(n\geq1\), relative to the endpoint and lower-index data.
\end{corollary}

\begin{proof}
For \(q\geq3\), the element \(-1\) is a nonidentity involution in
\(U(q)\). Character orthogonality gives
\(\sum_\chi\chi(-1)=0\), and each summand is \(1\) or \(-1\).
There are consequently \(\varphi(q)/2\) characters of each parity.
Reflection retains the parity stated above. The principal
character is even, so the endpoint removes one additional
dimension exactly when \(k\) is odd. This proves
\eqref{dist:eq:dist-parity-dimension}.

For the negative ladder, differentiation of the Hurwitz
multiplication formula at \(s=-k\) gives
\begin{equation}
 \sum_{j=0}^{d-1}
 \zeta'\!\left(-k,\frac{a+j}{d}\right)
 =d^{-k}\bigl(\zeta'(-k,a)+\log d\,\zeta(-k,a)\bigr).
 \label{dist:eq:dist-negative-jet}
\end{equation}
The Bernoulli-polynomial values \(\zeta(-k,a)\) are lower data.
Thus the homogeneous distribution weight is \(t_p=p^{-k}\).
The homogeneous reflection is
\eqref{dist:eq:dist-reflection-rows}; an explicit right-hand side is
given below. This applies the first part to the ladder in the
source article.

For the parameter derivative tower, the Laurent expansion is
\begin{equation}
 \zeta(s,a)=\frac1{s-1}
   +\sum_{n\geq0}\frac{(-1)^n}{n!}\gamma_n(a)(s-1)^n.
 \label{dist:eq:dist-stieltjes-laurent}
\end{equation}
For \(k\geq1\), its pole disappears after differentiation in \(a\):
\begin{equation}
 F_k(\epsilon,a):=
 \partial_a^k\zeta(1+\epsilon,a)
 =(-1)^k(1+\epsilon)_k\zeta(k+1+\epsilon,a)
 =\sum_{n\geq0}\frac{(-1)^n}{n!}
          \gamma_n^{(k)}(a)\epsilon^n.
 \label{dist:eq:dist-stieltjes-generating}
\end{equation}
This standard identity is the starting point of the all-order
parameter derivative formula of Coffey \cite{stieltjes:Coffey}.
Differentiating the Hurwitz multiplication formula \(k\) times
in \(a\), including its chain-rule factors, gives
\[
 \sum_{j=0}^{d-1}
 F_k\!\left(\epsilon,\frac{a+j}{d}\right)
 =d^{k+1+\epsilon}F_k(\epsilon,a).
\]
Comparing coefficients yields the exact triangular identity
\begin{equation}
 \sum_{j=0}^{d-1}
 \gamma_n^{(k)}\!\left(\frac{a+j}{d}\right)
 =d^{k+1}\sum_{r=0}^{n}
   \binom nr(-\log d)^{n-r}\gamma_r^{(k)}(a).
 \label{dist:eq:dist-stieltjes-multiplication}
\end{equation}
With lower indices \(r<n\) treated as specified data, the
homogeneous weight is \(t_p=p^{k+1}\), independently of \(n\).
Equation~\eqref{dist:eq:dist-anchor-dimension} now gives rank
\(q-\varphi(q)\).
\end{proof}

For completeness, the right-hand side of the negative reflection
can be written without introducing another Hurwitz derivative.
For \(0<a<1\), set
\[
 C_j(a)=\sum_{m\geq1}\frac{\cos(2\pi ma)}{m^j},
 \qquad
 S_j(a)=\sum_{m\geq1}\frac{\sin(2\pi ma)}{m^j}.
\]
For \(k\geq1\), these series at \(j=k+1\) converge absolutely, and
the Fourier form of the Hurwitz functional equation \cite{hstruct:DLMF}
gives
\begin{equation}
 \zeta'(-k,a)+(-1)^k\zeta'(-k,1-a)
 =\frac{(-1)^{\lfloor k/2\rfloor}k!}{(2\pi)^k}
 \begin{cases}
 C_{k+1}(a),&k\ \text{even},\\
 S_{k+1}(a),&k\ \text{odd}.
 \end{cases}
 \label{dist:eq:dist-negative-reflection}
\end{equation}
Indeed, the Fourier expansion of \(\zeta(s,a)\) has the factor
\[
 \frac{2\Gamma(1-s)}{(2\pi)^{1-s}}
 \left(\sin\frac{\pi s}{2}\,C_{1-s}(a)
       +\cos\frac{\pi s}{2}\,S_{1-s}(a)\right).
\]
The indicated reflection selects a sine or cosine prefactor
that vanishes at \(s=-k\). Only the derivative of that prefactor
survives, giving \eqref{dist:eq:dist-negative-reflection}.
This also explains why reflection is an inhomogeneous relation
with depth-one data in the inventory.

\begin{corollary}[The formal gamma grid]
\label{dist:cor:dist-gamma-rank}
For \(q\geq3\), the \(q-1\) formal coordinates
\(\log\Gamma(a/q)\), \(1\leq a<q\), modulo multiplication and
reflection and relative to their stated elementary right-hand
sides, have quotient dimension \(\varphi(q)/2\).
\end{corollary}

\begin{proof}
Lerch's identity
\(\zeta'(0,a)=\log\Gamma(a)-\tfrac12\log(2\pi)\)
translates gamma multiplication into
\eqref{dist:eq:dist-negative-jet} with \(k=0\), hence \(t_p=1\).
The reflection formula
\[
 \zeta'(0,a)+\zeta'(0,1-a)=-\log(2\sin\pi a)
\]
has homogeneous part \([x]+[-x]=0\). Apply
\eqref{dist:eq:dist-parity-dimension} with \(k=0\).
\end{proof}

\begin{remark}[What these rank theorems settle]
\label{dist:rem:dist-formal-only}
The first Stieltjes source states a finite verification followed
by an all-denominator conjecture; the parameter-derivative source
states a general rank theorem while reporting only a finite
calculation. Corollary~\ref{dist:cor:dist-all-ranks} supplies a proof
for the specified relation modules in both cases. The README
had already identified the second proof gap.
Corollary~\ref{dist:cor:dist-gamma-rank} likewise determines the rank
of the stated gamma relation matrix. None of these results says
that the evaluated numbers admit no further relations.
In particular, the claim that multiplication and reflection
generate every algebraic multiplicative relation among rational
gamma values remains part of the conjectural completeness
problem, not a consequence of the matrix rank.
\end{remark}

\subsection{Analytic weights, exceptional coordinates, and finite jets}

Specialize \(t_p=p^s=\exp(s\log p)\), using the real logarithm
of the prime. Formula~\eqref{dist:eq:dist-normal-coefficient} becomes
\begin{equation}
 A_{d,\chi}(s)=
 \left(\frac d{f_\chi}\right)^s
 \prod_{\substack{p\mid d\\p\nmid f_\chi}}
       \left(1-\chi^*(p)p^{-s}\right).
 \label{dist:eq:dist-euler-coefficient}
\end{equation}
This is exactly the familiar imprimitive Euler-factor pattern
for Dirichlet \(L\)-functions \cite{hstruct:DLMF}. The normal-form proof
shows why it appears in the formal relations before any
evaluation is made.

The symbols of exact denominator \(q\) have character coordinates
\(Y_{q,\chi}\). In the quotient, their \(\chi\)-coordinate is
\begin{equation}
 Y_{q,\chi}=A_{q,\chi}(s)Y_{f_\chi,\chi}.
 \label{dist:eq:dist-top-coordinate}
\end{equation}
They form a basis of the full quotient precisely when every
\(A_{q,\chi}(s)\) is nonzero. In particular this holds when
\(\Real s\ne0\), since \(\chi^*(p)\) has modulus one.
All nonzero integer weights in the two Stieltjes articles are
therefore free of this coordinate obstruction. On the imaginary
axis the top-denominator coordinates can fail, while the
conductor coordinates of Theorem~\ref{dist:thm:dist-normal-form}
remain a basis.

\begin{theorem}[Exact defect of top-denominator jets]
\label{dist:thm:dist-jet-defect}
Fix \(s_0\in\C\) and \(N\geq0\), and work over
\(\mathcal A_N=\C[\epsilon]/(\epsilon^{N+1})\) with
\(t_p=p^{s_0+\epsilon}\), truncated in this ring.
The full distribution module is free of rank \(\varphi(q)\)
over \(\mathcal A_N\), hence has complex dimension
\((N+1)\varphi(q)\).
Set
\begin{equation}
 r_\chi(s_0)=
 \#\{p\mid q:p\nmid f_\chi,\ p^{s_0}=\chi^*(p)\}.
 \label{dist:eq:dist-resonance-count}
\end{equation}
The image of the top-denominator submodule in its \(\chi\)-block
has complex dimension
\begin{equation}
 \max\{N+1-r_\chi(s_0),0\},
 \label{dist:eq:dist-jet-image}
\end{equation}
and its cokernel has dimension
\(\min\{r_\chi(s_0),N+1\}\).
An endpoint anchor or a fixed-parity reflection restricts these
formulas to the retained characters. The resulting full quotient
remains free over \(\mathcal A_N\) with the ranks already computed.
\end{theorem}

\begin{proof}
Freeness is Theorem~\ref{dist:thm:dist-normal-form}, applied to
\(\mathcal A_N\). A resonant factor in
\eqref{dist:eq:dist-euler-coefficient} has a simple zero at \(s_0\):
\[
 \left.\frac{d}{ds}
 (1-\chi^*(p)p^{-s})\right|_{s=s_0}=\log p\ne0.
\]
All other factors, including \((q/f_\chi)^s\), are nonzero
there. In the ring of analytic germs,
\[
 A_{q,\chi}(s_0+\epsilon)
   =\epsilon^{r_\chi(s_0)}u_\chi(\epsilon),
 \qquad u_\chi(0)\ne0.
\]
After truncation, \(u_\chi\) is a unit. The
top-denominator image in the free rank-one block is therefore
the ideal \(\epsilon^{r_\chi(s_0)}\mathcal A_N\).
Its basis consists of the powers
\(\epsilon^{r_\chi},\ldots,\epsilon^N\) when \(r_\chi\leq N\),
and it is zero otherwise. This proves both dimensions.
Anchoring and reflection remove entire character blocks,
so the same calculation applies to each retained block.
\end{proof}

Comparing Taylor coefficients of a weighted distribution identity
gives the corresponding stacked finite-jet system. If one identifies
the coefficient at order \(j\) with
\(\epsilon^{N-j}[x]\), its relation space is the
\(\mathcal A_N\)-span above; factorial normalizations merely
rescale coordinates. Thus freeness also gives the rank of the
complete simultaneous jet system, rather than just its successive
highest layers.

\begin{proposition}[Multiplicity of nonzero resonances]
\label{dist:prop:dist-resonance-multiplicity}
Resonances occur only on \(\Real s_0=0\).
For \(s_0\ne0\), each \(r_\chi(s_0)\) is at most one.
Multiple resonant primes in a single character block can occur
only at \(s_0=0\).
\end{proposition}

\begin{proof}
Taking absolute values in \(p^{s_0}=\chi^*(p)\) gives
\(\Real s_0=0\). Write \(s_0=\iu t\), with \(t\ne0\).
Because every character value is a root of unity, resonance
at \(p\) implies
\(t\log p/(2\pi)\in\Q\). This rational number is nonzero.
If distinct primes \(p\) and \(\ell\) both resonate, then
\(\log p/\log\ell\in\Q\). Writing this positive rational as
\(a/b\) would give \(p^b=\ell^a\) for positive integers \(a,b\),
contradicting unique factorization.
\end{proof}

For a concrete origin resonance, take \(q=14\) and let \(\chi^*\)
be the quadratic character modulo \(7\). Its conductor is \(7\),
its parity is odd, and \(\chi^*(2)=1\). Then
\begin{equation}
 Y_{14,\chi}(s)=(2^s-1)Y_{7,\chi}(s).
 \label{dist:eq:dist-resonance-example}
\end{equation}
At \(s=0\), the top-denominator coordinate vanishes although the
conductor direction persists. For any analytic realization of
this relation,
\begin{align}
 Y_{14,\chi}'(0)&=\log2\,Y_{7,\chi}(0),\notag\\
 Y_{14,\chi}''(0)&=
 2\log2\,Y_{7,\chi}'(0)+(\log2)^2Y_{7,\chi}(0).
 \label{dist:eq:dist-resonance-example-jets}
\end{align}
Thus recovering the first conductor derivative from these
top-denominator data requires the second top-denominator jet.
The actual realization
\[
 Y_{d,\chi}(s)
 =\sum_{a\in U(d)}\chi_d(a)\zeta(s,a/d)
 =d^sL(s,\chi_d)
\]
obeys the same identity. This example concerns an imprimitive
Euler factor; it is distinct from the trivial-zero regularization
of the primitive \(L\)-function in the next section.

\section{Primitive-grid determinants, local defects and finite descent}\label{jets:sec:det}
The polynomial character normal form above remains valid over arbitrary
complex coefficient algebras. Here identify the prime weights $A_p$ with
$t_p$, write $\mathcal D_q$ for that quotient, and put
$v_{m,\chi}=Y_{m,\chi}$, $b_\chi=Y_{f,\chi}$ for
$f=\operatorname{cond}(\chi)$. The coefficient $E_{m,\chi}$ is
$A_{m,\chi}$ of Theorem~\ref{dist:thm:dist-normal-form}. Thus
\begin{equation}\label{jets:eq:conductor-normal}
 [v_{m,\chi}]=E_{m,\chi}b_\chi.
\end{equation}
No Euler-factor division is used.
Let $U_q=\{u\bmod q:(u,q)=1\}$, and use $r=0$ for the Hurwitz endpoint
$x_r=1$. Distinguish the rank of the full quotient from the image of
its primitive-grid coordinate map.

Conductor coordinates are polynomial and globally valid, but it is often
more convenient to use the $\varphi(q)$ primitive-grid points
$e_{u/q}$, $u\in U_q$. Fourier-transforming those points gives
$v_{q,\chi}$. By Theorem~\ref{dist:thm:dist-normal-form}, the map from their formal
span into $\D_q$ is diagonal in character coordinates:
\begin{equation}\label{jets:eq:primitive-diagonal}
 v_{q,\chi}\longmapsto E_{q,\chi}b_\chi.
\end{equation}
Its determinant in these specified Fourier and conductor bases is
$\Delta_q=\prod_{\chi\bmod q}E_{q,\chi}$.

For $p^e\Vert q$, write
\[
 M_p=q/p^e,\qquad h_p=\ord_{M_p}(p),
\]
with $h_p=1$ when $M_p=1$.

\begin{theorem}[Determinant factorization]\label{jets:thm:determinant}
One has the polynomial identity
\begin{equation}\label{jets:eq:determinant}
 \boxed{\displaystyle
 \Delta_q(\mathbf A)=
 \prod_{p^e\Vert q}
 A_p^{\varphi(M_p)(p^{e-1}-1)}
 \bigl(A_p^{h_p}-1\bigr)^{\varphi(M_p)/h_p}.}
\end{equation}
Thus primitive-grid coordinates form a basis precisely when all the
factors in \eqref{jets:eq:determinant} are invertible in a field
specialization. Their failure does not change the rank of $\D_q$.
\end{theorem}
\begin{proof}
Fix $p^e\Vert q$ and inspect the $p$ factor of
$\prod_\chi E_{q,\chi}$. Characters whose conductor is not divisible by
$p$ are exactly the characters of $U_{M_p}$. Their additional factors
are $A_p-\chi(p)$. In the character group of $U_{M_p}$, the values
$\chi(p)$ run through every $h_p$th root of unity with multiplicity
$\varphi(M_p)/h_p$. This follows by duality from the fact that $p$ has
order $h_p$ in $U_{M_p}$. Their product is therefore
$(A_p^{h_p}-1)^{\varphi(M_p)/h_p}$.

It remains to count the pure powers of $A_p$. If $a=v_p(f)$, the
exponent contributed by one character is $e-\max(1,a)$. Writing this
as a sum of indicators gives
\[
 \sum_{\chi\bmod q}\bigl(e-\max(1,v_p(f_\chi))\bigr)
 =\sum_{j=1}^{e-1}\#\{\chi:v_p(f_\chi)\le j\}.
\]
The count in the last sum is $\varphi(p^j)\varphi(M_p)$, since these
are exactly the characters induced from level $p^jM_p$. Hence the
exponent is
$\varphi(M_p)\sum_{j=1}^{e-1}\varphi(p^j)
=\varphi(M_p)(p^{e-1}-1)$. Multiply over the primes.
\end{proof}

\begin{example}[Level twelve]
Let $\chi_3$ and $\chi_4$ be the nonprincipal quadratic characters of
conductors three and four, and let $\chi_{12}=\chi_3\chi_4$.
Their conductor generators are
\begin{align*}
 b_3&=e_{1/3}-e_{2/3},&
 b_4&=e_{1/4}-e_{3/4},\\
 b_{12}&=e_{1/12}-e_{5/12}-e_{7/12}+e_{11/12}.
\end{align*}
For $r\in U_{12}$, the normal form reads
\begin{align}\label{jets:eq:q12-normal}
 [e_{r/12}]=\frac14\bigl[&A_2(A_2-1)(A_3-1)e_0
 +\chi_3(r)A_2(A_2+1)b_3\notag\\
 &+\chi_4(r)(A_3+1)b_4+\chi_{12}(r)b_{12}\bigr].
\end{align}
Consequently
\[
 \Delta_{12}=A_2^2(A_2^2-1)(A_3^2-1).
\]
At $A_2=A_3=1$, the primitive-grid span misses the principal conductor
coordinate, even though the full quotient still has rank four.
\end{example}

\subsection{Local jet defects of the primitive-grid map}
Specialize to $A_p(s)=p^s$, and let $s=s_0+t$. All exponential factors
are nonzero. A factor $p^s-\chi(p)$ vanishes precisely when
$p^{s_0}=\chi(p)$, and any such zero is simple because its derivative
is $p^{s_0}\log p\ne0$.

\begin{theorem}[Local Smith factors and primitive-jet ranks]\label{jets:thm:smith}
Define
\begin{equation}\label{jets:eq:rho-general}
 \rho_\chi(s_0)=\#\{p\mid q:p\nmid f_\chi,\ p^{s_0}=\chi(p)\}.
\end{equation}
Over $\C[[t]]$, the diagonal factors of the primitive-grid map are
units times $t^{\rho_\chi(s_0)}$. Consequently its rank on order-$N$
truncated jets, as a complex linear map, is
\begin{equation}\label{jets:eq:primitive-jet-rank}
 \sum_{\chi\bmod q}\max\{N+1-\rho_\chi(s_0),0\}.
\end{equation}
At $s_0=0$, this is determined by
$\rho_\chi(0)=\#\{p\mid q:p\nmid f_\chi,\ \chi(p)=1\}$.
\end{theorem}
\begin{proof}
Factor $E_{q,\chi}(s_0+t)$ into its simple vanishing factors and a unit.
The primitive Fourier matrix is constant and invertible, so these are
also the local invariant factors, up to permutation and units.
Multiplication by $t^r$ on $\C[t]/(t^{N+1})$ has rank
$\max(N+1-r,0)$. Summing over the diagonal coordinates proves the
formula. This is the rank of a \emph{coordinate inclusion map}, not
the rank of the distribution relation matrix.
\end{proof}

\begin{corollary}[Multiple-prime resonances are central]\label{jets:cor:central}
If $s_0\ne0$, at most one prime can satisfy
$p^{s_0h_p}=1$. At a nonzero resonant point for that prime $p$, the
primitive-grid rank defect is $\varphi(M_p)/h_p$, and every affected
local factor has order one. The order-$N$ inclusion rank is then
$(N+1)\varphi(q)-\varphi(M_p)/h_p$ for all $N\ge0$.
\end{corollary}
\begin{proof}
A resonance forces $s_0$ to be purely imaginary. Resonances at distinct
primes would make $\log p/\log\ell$ rational, because each phase is a
rational multiple of $2\pi$. Unique prime factorization rules this out
unless $s_0=0$. For the one resonant prime, the multiplicity count in
Theorem~\ref{jets:thm:determinant} gives the number of affected characters.
Apply Theorem~\ref{jets:thm:smith}.
\end{proof}

\begin{example}[Two exact jet-rank profiles]
At level $12$ and $s_0=0$, the principal character has $\rho=2$ and the
other three characters have $\rho=0$. For $N=0,1,2$, the primitive-jet
ranks are $3,6,10$, while the full quotient dimensions are $4,8,12$.
At level $21$, the principal character has $\rho=2$, the character
induced from conductor three has $\rho=1$, and the other ten
characters have $\rho=0$. The inclusion ranks are $10,21,33$, compared
with full quotient dimensions $12,24,36$. Indeed,
\[
 \Delta_{21}=(A_3^6-1)(A_7-1)^2.
\]
Exact character-phase certificates for these examples are supplied.
\end{example}

\section{Explicit primitive-grid reduction coefficients}\label{jets:sec:primitive}

Away from the exceptional locus, the conductor normal form can be
inverted to express all grid points through the primitive ones. For
Hurwitz weights there is a particularly concrete answer that avoids
character-field arithmetic.

Let $\mathcal S_q$ be the multiplicative semigroup of positive integers
whose prime factors divide $q$, including $1$. For $\Real s>0$, define
\begin{equation}\label{jets:eq:C-smooth}
 C_{r,u}(s)=\sum_{\substack{d\in\mathcal S_q\\du\equiv r\; (q)}}d^{-s},
 \qquad 0\le r<q,\quad u\in U_q.
\end{equation}
The series is absolutely convergent, since
$\sum_{d\in\mathcal S_q}d^{-\sigma}
=\prod_{p\mid q}(1-p^{-\sigma})^{-1}$ for $\sigma>0$.

\begin{theorem}[Primitive-grid extension]\label{jets:thm:primitive}
For $\Real s>0$, the $q\times\varphi(q)$ matrix $C(s)$ satisfies every
prime and divisor distribution relation with weights $p^s$, and its
primitive rows form the identity matrix. It is the unique extension
matrix from primitive-grid data. In particular,
\begin{equation}\label{jets:eq:Hurwitz-C}
 \zeta(s,r/q)=\sum_{u\in U_q}C_{r,u}(s)\zeta(s,u/q)
\end{equation}
with the endpoint convention for $r=0$, initially for $\Real s>1$ and
then meromorphically wherever the coefficients are defined.
\end{theorem}
\begin{proof}
For a primitive row $r$, the congruence $du\equiv r\pmod q$ forces
$d$ to be prime to $q$, hence $d=1$. This gives the identity submatrix.
For a prime distribution row indexed by $a\equiv0\pmod p$,
\[
 \sum_{pr\equiv a\;(q)}C_{r,u}(s)
 =\sum_{\substack{d\in\mathcal S_q\\pdu\equiv a\;(q)}}d^{-s}
 =p^s\sum_{\substack{d'\in\mathcal S_q\\d'u\equiv a\;(q)}}(d')^{-s}.
\]
In the last step $d'=pd$; conversely every $d'$ in the last sum is
divisible by $p$, because $u$ is a unit and $a$ is divisible by $p$.
The last expression is $p^sC_{a,u}(s)$. Prime sufficiency gives all
divisor rows. Theorem~\ref{dist:thm:dist-normal-form} and the identity submatrix
show that these $\varphi(q)$ independent columns exhaust the solution
space. Hurwitz zeta satisfies the same distribution equations by
regrouping its absolutely convergent defining series, proving
\eqref{jets:eq:Hurwitz-C}.
\end{proof}

\subsection{A finite formula, including negative spectral parameters}
For $0\le r<q$, let
\[
 g=\gcd(r,q),\quad m=q/g,\quad a=r/g,
 \quad P_r=\{p\mid q:p\nmid m\}.
\]
When $r=0$, this means $g=q$, $m=1$, $a=0$.
For $p\in P_r$, put $h_{p,m}=\ord_m(p)$, with order one modulo one.

\begin{theorem}[Finite residue formula]\label{jets:thm:finite-C}
The meromorphic continuation of \eqref{jets:eq:C-smooth} is
\begin{equation}\label{jets:eq:C-finite}
 C_{r,u}(s)=g^{-s}
 \frac{\displaystyle
 \sum_{\substack{0\le b_p<h_{p,m}\;(p\in P_r)\\
 u\prod_{p\in P_r}p^{b_p}\equiv a\;(m)}}
 \prod_{p\in P_r}p^{-s b_p}}
 {\displaystyle\prod_{p\in P_r}(1-p^{-s h_{p,m}})}.
\end{equation}
Empty products have their usual value one. In particular,
\begin{equation}\label{jets:eq:C-endpoint}
 C_{0,u}(s)=q^{-s}\prod_{p\mid q}(1-p^{-s})^{-1},
\end{equation}
independently of $u$.
\end{theorem}
\begin{proof}
The congruence in \eqref{jets:eq:C-smooth} forces the valuation of $d$ at
an unsaturated prime, one dividing $m$, to be exactly $v_p(g)$.
At a saturated prime, one not dividing $m$, it is at least $v_p(g)$.
Thus $d=g\prod_{p\in P_r}p^{c_p}$ with $c_p\ge0$, and the remaining
congruence is $u\prod p^{c_p}\equiv a\pmod m$.
Write $c_p=b_p+h_{p,m}j_p$ and sum the independent geometric series in
$j_p$. This proves the formula in $\Real s>0$ and hence its meromorphic
continuation. The endpoint has trivial congruence modulo one.
\end{proof}

Formula \eqref{jets:eq:C-finite}, \emph{not} the divergent smooth-number
series, is the definition used at negative real $s$. It has no pole at
any nonzero real $s$, and the continued distribution identities remain
valid there. At a zero of a displayed denominator, one should use the
polynomial conductor normal form rather than cancel poles numerically.
Some individual coefficients may have removable singularities, but
Theorem~\ref{jets:thm:determinant} gives the criterion for the entire
primitive-grid coordinate system.

\begin{corollary}[Exact support and signs]\label{jets:cor:support}
Let $H_r$ be the subgroup of $U_m$ generated by the classes of the
primes in $P_r$. For every nonzero real $s$,
\[
 C_{r,u}(s)\ne0\quad\Longleftrightarrow\quad au^{-1}\in H_r.
\]
The number of supported primitive coordinates is
\begin{equation}\label{jets:eq:support-count}
 \frac{\varphi(q)}{\varphi(m)}|H_r|.
\end{equation}
All supported coefficients are positive for $s>0$. For $s<0$, they
all have sign $(-1)^{|P_r|}$.
\end{corollary}
\begin{proof}
The numerator of \eqref{jets:eq:C-finite} is a sum of positive real terms
and is nonempty exactly under the subgroup condition. The reduction
map $U_q\to U_m$ has fibers of cardinality
$\varphi(q)/\varphi(m)$, giving \eqref{jets:eq:support-count}.
For $s>0$ every denominator factor is positive. For $s<0$ each of the
$|P_r|$ denominator factors is negative, while the numerator and
$g^{-s}$ are positive.
\end{proof}

\begin{corollary}[Complete monotonicity and normalization]\label{jets:cor:monotone}
For $s>0$ and $j\ge0$, put
\begin{equation}\label{jets:eq:W-definition}
 W_{r,u}^{[j]}(s)=(-1)^j\frac{d^j}{ds^j}C_{r,u}(s)
 =\sum_{\substack{d\in\mathcal S_q\\du\equiv r\;(q)}}
 \frac{(\log d)^j}{d^s}.
\end{equation}
These quantities are nonnegative. Moreover,
$\sum_u C_{r,u}(1)=1$ for every $r$, and
$0\le C_{r,u}(k)\le1$ for integer $k\ge1$.
\end{corollary}
\begin{proof}
The logarithmically weighted semigroup series converges locally
uniformly in $\Real s>0$, so termwise differentiation is valid.
The constant grid vector satisfies the distribution equations at
$s=1$, and its primitive entries are all one. Uniqueness of extension
proves the row-sum identity. Monotonicity gives
$C_{r,u}(k)\le C_{r,u}(1)\le1$.
\end{proof}

\begin{example}[Two reductions at denominator six]\label{jets:ex:q6}
The nontrivial coefficients for $r=2$ are
\[
 C_{2,1}(s)=\frac{2^s}{2^{2s}-1},\qquad
 C_{2,5}(s)=\frac1{2^{2s}-1}.
\]
Thus, meromorphically,
\begin{align}
 \zeta(s,1/3)&=
 \frac{2^s\zeta(s,1/6)+\zeta(s,5/6)}{2^{2s}-1},\label{jets:eq:q6-one}\\
 \zeta(s,1/2)&=
 \frac{\zeta(s,1/6)+\zeta(s,5/6)}{3^s-1}.
 \label{jets:eq:q6-two}
\end{align}
At $s=1$, the coefficients in \eqref{jets:eq:q6-one} are $2/3$ and $1/3$.
At $s=-1$ they are $-2/3$ and $-4/3$, illustrating the sign theorem
without using a divergent series. The poles at $s=0$ reflect a
singular primitive basis, not a failure of the underlying distribution
relations.
\end{example}

\subsection{Denominator certificates and construction cost}
Let $M_p$ also denote, in this paragraph only, the pushforward matrix
on grid vectors defined by
\[
 (M_pf)_r=\sum_{pu\equiv r\;(q)}f_u.
\]
To avoid confusion with the integer complement in
Section~\ref{jets:sec:det}, write that complement as $q_p=q/p^e$ here and
put $h_p=\ord_{q_p}(p)$. Multiplication of residues gives
$M_p^{e+h_p}=M_p^e$. For $z=p^{-s}$,
\begin{equation}\label{jets:eq:operator-inverse}
 (I-zM_p)^{-1}
 =\sum_{j=0}^{e-1}z^jM_p^j
 +\frac{z^e}{1-z^{h_p}}\sum_{j=0}^{h_p-1}z^jM_p^{e+j}.
\end{equation}
This follows first from a geometric series for $|z|<1$, then as a
rational matrix identity. The commuting product of these inverses,
applied to vectors supported on $U_q$, is the extension $C(s)$.

\begin{proposition}[A uniform rational denominator]\label{jets:prop:denominator}
For integer $k\ge1$, every reduced denominator of $C_{r,u}(k)$ divides
\begin{equation}\label{jets:eq:denominator-bound}
 D_q(k)=\prod_{p^e\Vert q}
 p^{k(e-1)}\bigl(p^{kh_p}-1\bigr).
\end{equation}
Its bit length is $O(kq\log q)$. The coefficients of the raw derivative
$C_{r,u}^{(j)}(k)$, viewed as a polynomial in the symbols $\log p$,
have denominators dividing $D_q(k)^{j+1}$.
\end{proposition}
\begin{proof}
Every scalar coefficient on the right of
\eqref{jets:eq:operator-inverse} has denominator dividing
$p^{k(e-1)}(p^{kh_p}-1)$. Multiplication over primes proves the first
claim. Since $h_p\le q$ and $\sum_{p\mid q}\log p\le\log q$, its
logarithm is at most a constant times $kq\log q$.
The matrix product can be written as a ratio of integer polynomials in
$p^s$, with denominator obtained from \eqref{jets:eq:denominator-bound} by
replacing $k$ with $s$. Repeated differentiation introduces polynomial
factors in $\log p$ and raises the denominator power by at most one
each time. For normalized Taylor coefficients, an additional factor
$j!$ may be needed.
\end{proof}

The implementation computes \eqref{jets:eq:C-finite} by convolution on
residue classes modulo $m$, separately for each $m\mid q$. At a fixed
$m$, each prime update takes at most $\varphi(m)^2$ rational arithmetic
operations. A complete matrix can therefore be constructed in
\[
 O\!\left(\omega(q)\sum_{m\mid q}\varphi(m)^2+q\varphi(q)\right)
 =O\bigl(\omega(q)q^2\bigr)
\]
rational arithmetic operations, using $\sum_{m\mid q}\varphi(m)=q$.
This is an arithmetic-operation bound, not a constant-cost bit model
or a claim of optimality. The denominator estimate supplies a
conservative size control at positive integer orders.

\section{All-orders Hurwitz--Stieltjes identities}\label{jets:sec:stieltjes}

Write the Laurent expansion as
\begin{equation}\label{jets:eq:stieltjes-laurent}
 \zeta(1+t,x)=\frac1t+
 \sum_{n\ge0}\frac{(-1)^n}{n!}\gamma_n(x)t^n,
 \qquad x>0.
\end{equation}
Here $\gamma_n^{(k)}$ always means $k$ derivatives in $x$; a derivative
on $C(s)$ is in the spectral variable. The coefficient identities in
this section are analytic theorems, distinct from formal rank claims.

\begin{theorem}[Derived Stieltjes reduction]\label{jets:thm:derived-stieltjes}
For $k\ge1$, $n\ge0$, and $x_r=r/q$ with the endpoint convention,
\begin{equation}\label{jets:eq:derived-stieltjes}
 \boxed{\displaystyle
 \gamma_n^{(k)}(x_r)=
 \sum_{u\in U_q}\sum_{j=0}^n\binom nj
 W_{r,u}^{[j]}(k+1)\,
 \gamma_{n-j}^{(k)}(u/q).}
\end{equation}
All the coefficient functions are given by the finite formula
\eqref{jets:eq:C-finite} and its derivatives.
\end{theorem}
\begin{proof}
Differentiating Hurwitz zeta in its parameter yields
\[
 \partial_x^k\zeta(1+t,x)=(-1)^k(1+t)_k\zeta(k+1+t,x).
\]
The pole in \eqref{jets:eq:stieltjes-laurent} is independent of $x$, so for
$k\ge1$ the left side is
$\sum_n(-1)^n\gamma_n^{(k)}(x)t^n/n!$.
Apply \eqref{jets:eq:Hurwitz-C} at spectral parameter $k+1+t$ and multiply
by the common scalar $(-1)^k(1+t)_k$. Taylor multiplication gives
\eqref{jets:eq:derived-stieltjes}, because
$(-1)^j C_{r,u}^{(j)}=W_{r,u}^{[j]}$.
\end{proof}

This proof also explains why the harmonic-number master identity in
Chapter 8 is compatible with the same reduction matrix at every
Stieltjes order: the rising factorial is common to every rational
argument. It does not create new independent distribution coordinates.

\begin{theorem}[Underived Stieltjes reduction, with the pole term]\label{jets:thm:underived}
For $n\ge0$,
\begin{equation}\label{jets:eq:underived-stieltjes}
 \boxed{\displaystyle
 \gamma_n(x_r)=
 \sum_{u\in U_q}\sum_{j=0}^n\binom nj
 W_{r,u}^{[j]}(1)\gamma_{n-j}(u/q)
 -\frac1{n+1}\sum_{u\in U_q}W_{r,u}^{[n+1]}(1).}
\end{equation}
\end{theorem}
\begin{proof}
Apply \eqref{jets:eq:Hurwitz-C} at $s=1+t$. The coefficients are analytic
there. Multiplying their Taylor expansions by
\eqref{jets:eq:stieltjes-laurent}, the coefficient of $t^n$ receives the
usual product terms and also
$\sum_u C_{r,u}^{(n+1)}(1)/(n+1)!$ from the pole $1/t$.
Multiplication by $(-1)^n n!$ converts that additional term to the
negative last term of \eqref{jets:eq:underived-stieltjes}. The coefficient
of $1/t$ agrees because $\sum_u C_{r,u}(1)=1$.
\end{proof}

Omitting the last term would already give an incorrect formula for
$\gamma_0$. The point is not a rank correction: a known scalar right
side does not change a homogeneous rank. It is a correction to the
actual analytic value identity.

\begin{example}[An explicit first-order identity]
From \eqref{jets:eq:q6-one}, direct differentiation at $s=1$ gives
\[
 W_{2,1}^{[1]}(1)=\frac{10}{9}\log2,\qquad
 W_{2,5}^{[1]}(1)=\frac89\log2.
\]
Consequently
\begin{equation}\label{jets:eq:gamma0-six}
 \gamma_0(1/3)=\frac23\gamma_0(1/6)+\frac13\gamma_0(5/6)-2\log2.
\end{equation}
At $s=2$, the analogous coefficients give
\begin{align}\label{jets:eq:gamma1prime-six}
 \gamma_1'(1/3)={}&\frac4{15}\gamma_1'(1/6)
 +\frac1{15}\gamma_1'(5/6)\notag\\
 &+\frac{68\log2}{225}\gamma_0'(1/6)
 +\frac{32\log2}{225}\gamma_0'(5/6).
\end{align}
Since $\gamma_0'=-\psi'$, the signs can also be checked against the
parameter-differentiated digamma multiplication identity.
The verification program checks \eqref{jets:eq:underived-stieltjes} through
$n=2$ and \eqref{jets:eq:gamma1prime-six} by numerical parameter
differentiation of Stieltjes functions.
\end{example}

\section{Finite conductor descent for all Stieltjes layers}

For a primitive character $\chi$ of conductor $f\mid d$, define
\begin{equation}
G_{d,\chi}^{(n,k)}=
\sum_{a\in\U_d}\chi(a)\gamma_n^{(k)}(a/d).
\end{equation}
Set
\begin{equation}\label{conductor:eq:Rdf}
R(d,f)=\prod_{\substack{p\mid d\\p\nmid f}}p,
\qquad m_r=\frac{d}{fr}\quad(r\mid R(d,f)).
\end{equation}
Each $m_r$ is a positive integer. Upon substituting $t_p=p^s$,
\eqref{dist:eq:dist-normal-coefficient} becomes the entire exponential polynomial
\begin{align}\label{conductor:eq:Aanalytic}
A_{d,f,\chi}(s)
&=\left(\frac df\right)^s
\prod_{\substack{p\mid d\\p\nmid f}}(1-\chi(p)p^{-s})\notag\\
&=\sum_{r\mid R(d,f)}\mu(r)\chi(r)m_r^s.
\end{align}
In particular,
\begin{equation}\label{conductor:eq:Aderivative}
A_{d,f,\chi}^{(j)}(s)
=\sum_{r\mid R(d,f)}\mu(r)\chi(r)m_r^s\log^j m_r.
\end{equation}

\begin{theorem}[All-index, all-parameter-order descent]\label{conductor:thm:stieltjes}
For $n\ge0$, $k\ge1$ and $f\mid d$,
\begin{align}\label{conductor:eq:stieldescent}
G_{d,\chi}^{(n,k)}
={}&\sum_{r\mid R(d,f)}\mu(r)\chi(r)m_r^{k+1}\notag\\
&\quad\cdot\sum_{j=0}^{n}(-1)^j\binom nj\log^j m_r\,
G_{f,\chi}^{(n-j,k)}.
\end{align}
No Euler-factor denominator occurs. Equivalently,
\begin{equation}\label{conductor:eq:stielA}
G_{d,\chi}^{(n,k)}=
\sum_{j=0}^{n}(-1)^j\binom nj A_{d,f,\chi}^{(j)}(k+1)
G_{f,\chi}^{(n-j,k)}.
\end{equation}
\end{theorem}
\begin{proof}
For $\Real s>1$, sorting a Dirichlet series by residue classes gives
\[
\sum_{a\in\U_d}\chi(a)\zeta(s,a/d)=d^sL(s,\chi_d),
\]
where $\chi_d$ is the induced character modulo $d$. Removing the extra
Euler factors from its primitive $L$-series yields
\[
L(s,\chi_d)=L(s,\chi)
\prod_{\substack{p\mid d\\p\nmid f}}(1-\chi(p)p^{-s}).
\]
Thus the character sum at $d$ is $A_{d,f,\chi}(s)$ times the character
sum at $f$. This is also the analytic realization of
Theorem~\ref{dist:thm:dist-normal-form}.

Insert $s=k+1+\eps$, multiply by the common factor
$(-1)^k(1+\eps)_k$, and use \eqref{tower:eq:masterN}. Comparison of
$(-1)^n\eps^n/n!$ gives \eqref{conductor:eq:stielA}; the sign $(-1)^j$ comes
from the Stieltjes Laurent convention. Finally apply
\eqref{conductor:eq:Aderivative}.
\end{proof}

\subsection{Concrete second-index identities at level twelve}

The multipliers are
\[
A_{12,3,\chi_3}(s)=4^s+2^s,\qquad
A_{12,4,\chi_4}(s)=3^s+1.
\]
Therefore, at $n=2$, $k=1$,
\begin{align}\label{conductor:eq:examplegamma}
G_{12,\chi_3}^{(2,1)}
&=20G_{3,\chi_3}^{(2,1)}
-72\log2\,G_{3,\chi_3}^{(1,1)}
+68\log^22\,G_{3,\chi_3}^{(0,1)},\\
G_{12,\chi_4}^{(2,1)}
&=10G_{4,\chi_4}^{(2,1)}
-18\log3\,G_{4,\chi_4}^{(1,1)}
+9\log^23\,G_{4,\chi_4}^{(0,1)}.
\end{align}
For instance,
\[
G_{12,\chi_3}^{(2,1)}=
\gamma_2'(1/12)-\gamma_2'(5/12)+\gamma_2'(7/12)-\gamma_2'(11/12),
\]
and $G_{3,\chi_3}^{(2,1)}=\gamma_2'(1/3)-\gamma_2'(2/3)$.
These are exact identities between explicitly specified signed sums; they
require no integer-relation search. The two identities were also checked
numerically using independent finite Hurwitz character sums.

\subsection{The undifferentiated layer and its pole correction}

For $k=0$, a principal character sum still contains the Hurwitz pole.
The same formula without a correction would be wrong.

\begin{proposition}[The $k=0$ correction]\label{conductor:prop:kzero}
For nonprincipal primitive $\chi$ the formula
\eqref{conductor:eq:stielA} remains valid with $k=0$. For $f=1$, put
$A_d(s)=A_{d,1,1}(s)$. Then
\begin{equation}\label{conductor:eq:kzero}
G_{d,1}^{(n,0)}=
\sum_{j=0}^{n}(-1)^j\binom nj A_d^{(j)}(1)\gamma_{n-j}(1)
+\frac{(-1)^n}{n+1}A_d^{(n+1)}(1).
\end{equation}
\end{proposition}
\begin{proof}
For a nonprincipal character, the sum of its values over units is zero,
so the $1/\eps$ residues cancel. Coefficient comparison is unchanged.
For the principal conductor, multiply
\[
A_d(1+\eps)\left(\frac1\eps+
\sum_{n\ge0}\frac{(-1)^n\gamma_n(1)}{n!}\eps^n\right).
\]
The pole term contributes $A_d^{(n+1)}(1)/(n+1)!$ to the coefficient of
$\eps^n$. Multiplying by $(-1)^n n!$ gives the correction in
\eqref{conductor:eq:kzero}. The residue itself is
$A_d(1)=d\prod_{p\mid d}(1-p^{-1})=\varphi(d)$, as required.
\end{proof}

For the non-character multiplication formula this specializes to
\begin{equation}\label{conductor:eq:classicalstieldist}
\sum_{j=0}^{d-1}\gamma_n\!\left(\frac{a+j}{d}\right)
=d\sum_{r=0}^{n}\binom nr(-\log d)^r\gamma_{n-r}(a)
+\frac{(-1)^n d\log^{n+1}d}{n+1}.
\end{equation}
The inhomogeneous term is the residue contribution, not an optional
normalization.


\section{Cyclotomic polylogarithm traces}\label{jets:sec:traces}

The same finite module applies to a second evaluation. Put
$\zeta_m=\exp(2\pi\iu/m)$ and evaluate
$e_x\mapsto\Li_s(\ee^{2\pi\iu x})$.
For $\Real s>1$, grouping the defining series by multiples of $p$ gives
\[
 \sum_{py=x}\Li_s(\ee^{2\pi\iu y})
 =p^{1-s}\Li_s(\ee^{2\pi\iu x}).
\]
Thus the weights are now $A_p=p^{1-s}$, and analytic continuation
extends the identity wherever the individual expressions are defined.

For a primitive character $\chi$ of conductor $f\mid q$, define
\begin{equation}\label{jets:eq:trace-definition}
 \cP_{q,\chi}(s)=
 \sum_{u\in U_q}\chi(u)\Li_s(\zeta_q^u),\qquad
 \tau(\chi)=\sum_{a\in U_f}\chi(a)\zeta_f^a.
\end{equation}
At conductor one take $\tau(\mathbf1)=1$ and $L(s,\mathbf1)=\zeta(s)$.

\begin{theorem}[Conductor-lifted trace identity]\label{jets:thm:trace}
As a meromorphic identity, and with removable singularities interpreted
by continuation,
\begin{equation}\label{jets:eq:trace-formula}
 \boxed{\displaystyle
 \cP_{q,\chi}(s)=
 (q/f)^{1-s}
 \prod_{\substack{p\mid q\\p\nmid f}}
 \bigl(1-\chi(p)p^{s-1}\bigr)
 \tau(\chi)L(s,\overline\chi).}
\end{equation}
For every $q>1$, the left side is an entire function of the spectral
variable $s$.
\end{theorem}
\begin{proof}
At primitive conductor $f>1$, absolute convergence and the primitive
Gauss-sum identity give
\[
 \sum_{a\in U_f}\chi(a)\Li_s(\zeta_f^a)
 =\sum_{n\ge1}\frac{\sum_a\chi(a)\zeta_f^{an}}{n^s}
 =\tau(\chi)L(s,\overline\chi).
\]
The Gauss-sum convention, including the complex conjugate on $\chi$,
is the one in~\cite{jets:dlmf-dirichlet}. Applying
\eqref{jets:eq:conductor-normal} with $A_p=p^{1-s}$ multiplies this
primitive trace by $E_{q,\chi}$. Extracting $p^{1-s}$ from each
first-appearance factor gives precisely the factor in
\eqref{jets:eq:trace-formula}. The conductor-one case follows directly
from $\Li_s(1)=\zeta(s)$ and the same polynomial identity.

For completeness, when $z=\zeta_q^u\ne1$, regrouping by residues gives
\begin{equation}\label{jets:eq:polylog-Hurwitz-DFT}
 \Li_s(z)=q^{-s}\sum_{a=1}^q z^a\zeta(s,a/q).
\end{equation}
Hurwitz zeta has only its simple pole at $s=1$, whose coefficient in
this finite sum is $\sum_a z^a=0$. Hence $\Li_s(z)$ is entire in $s$.
All arguments in the primitive level-$q$ trace differ from one.
\end{proof}

\begin{remark}[The character phase is consequential]
For the order-four character modulo five with $\chi(2)=\iu$, the
level-ten lifting factor is $2^{1-s}-\iu$, not
$2^{1-s}+\iu$. A complex-character numerical check is included to
catch an accidental interchange of $\chi$ and $\overline\chi$.
Real quadratic examples alone would not detect that error.
\end{remark}

\subsection{Exact nonprincipal vanishing orders}
For a nonprincipal primitive $\chi$, set
\[
 \rho_\chi=\#\{p\mid q:p\nmid f,\ \chi(p)=1\}.
\]
These are the same conductor statistics that controlled the
primitive-grid jet defects at Hurwitz order zero.

\begin{theorem}[Character trace zeros at order one]\label{jets:thm:trace-zero}
The entire function $\cP_{q,\chi}(s)$ has a zero of exact order
$\rho_\chi$ at $s=1$, where order zero means a nonzero value.
All lower derivatives vanish, and
\begin{align}\label{jets:eq:trace-leading}
 \cP_{q,\chi}^{(\rho_\chi)}(1)
 ={}&(-1)^{\rho_\chi}\rho_\chi!\,
 \tau(\chi)L(1,\overline\chi)\notag\\
 &\times\prod_{\substack{p\mid q,\ p\nmid f\\\chi(p)=1}}\log p
 \prod_{\substack{p\mid q,\ p\nmid f\\\chi(p)\ne1}}(1-\chi(p)).
\end{align}
\end{theorem}
\begin{proof}
Each factor with $\chi(p)=1$ in \eqref{jets:eq:trace-formula} is
$1-p^{s-1}=-(s-1)\log p+O((s-1)^2)$.
All other displayed factors are nonzero at $s=1$.
A primitive Gauss sum is nonzero, and Dirichlet's nonvanishing theorem
states $L(1,\overline\chi)\ne0$ for nonprincipal characters
\cite[Eq.~25.15.9]{jets:dlmf-dirichlet}. Thus there is no additional zero.
Taking the first nonzero Taylor coefficient proves the formula.
\end{proof}

Increasing the exponents of primes already present in $q$ changes
higher Taylor coefficients but not this first nonzero derivative:
$(q/f)^{1-s}$ equals one at $s=1$. This invariance provides an
additional family of identities between traces at different levels.

\subsection{The principal trace and its pole cancellation}
Write $\cP_q=\cP_{q,\mathbf1}$ and define
\[
 r=\omega(q),\quad P=\prod_{p\mid q}\log p,\qquad
 B=\log q-\tfrac12\log\rad(q).
\]
The conductor-one factor contains a pole, so the vanishing order is
one less than the number of primes, not equal to that number.

\begin{theorem}[Principal trace jet generator]\label{jets:thm:principal}
In a neighborhood of $t=0$,
\begin{equation}\label{jets:eq:principal-generator}
 \boxed{\displaystyle
 \cP_q(1+t)=(-1)^rP\,t^{r-1}\,
 \ee^{-Bt}\,[t\zeta(1+t)]
 \prod_{p\mid q}\frac{\sinh(t\log p/2)}{t\log p/2}.}
\end{equation}
In particular its zero at $s=1$ has exact order $r-1$, and
\begin{align}
 \cP_q^{(j)}(1)&=0 &&(0\le j<r-1),\label{jets:eq:principal-vanish}\\
 \cP_q^{(r-1)}(1)&=(-1)^r(r-1)!P,\label{jets:eq:principal-first}\\
 \cP_q^{(r)}(1)&=(-1)^r r!P(\gamma-B),\label{jets:eq:principal-second}\\
 \cP_q^{(r+1)}(1)&=(-1)^r(r+1)!P
 \left[-\gamma_1-B\gamma+\frac{B^2}{2}
 +\frac1{24}\sum_{p\mid q}(\log p)^2\right].
 \label{jets:eq:principal-third}
\end{align}
Here $\gamma_1$ is the ordinary first Stieltjes constant in
\eqref{jets:eq:stieltjes-laurent}.
\end{theorem}
\begin{proof}
The principal case of \eqref{jets:eq:trace-formula} is
\[
 \cP_q(1+t)=q^{-t}\prod_{p\mid q}(1-p^t)\zeta(1+t).
\]
Use $1-\ee^x=-x\ee^{x/2}\sinh(x/2)/(x/2)$ in every factor.
This proves \eqref{jets:eq:principal-generator}.
The factors in square brackets and in the hyperbolic-sine product are
analytic and equal to one at zero. Expanding them gives
\[
 t\zeta(1+t)=1+\gamma t-\gamma_1t^2+O(t^3),\qquad
 \prod_p\frac{\sinh(t\log p/2)}{t\log p/2}
 =1+\frac{t^2}{24}\sum_p(\log p)^2+O(t^4).
\]
Multiplication by $\ee^{-Bt}$ gives all the displayed derivatives.
\end{proof}

\subsection{Concrete exact identities}
The following examples are proved consequences of the preceding
theorems, not integer-relation conjectures. The derivatives are in the
\emph{order} of the polylogarithm, not its argument.

At level thirty,
\begin{align}
 \sum_{u\in U_{30}}\Li_1(\zeta_{30}^u)&=0,\label{jets:eq:q30-zero}\\
 \sum_{u\in U_{30}}\Liord{1}{\zeta_{30}^u}&=0,\label{jets:eq:q30-first}\\
 \sum_{u\in U_{30}}\Liord{2}{\zeta_{30}^u}
 &=\boxed{-2\log2\log3\log5}.\label{jets:eq:q30-second}
\end{align}
The third derivative of this trace is
$-6\log2\log3\log5\,[\gamma-\frac12\log30]$.
At level $210$, derivatives of orders zero, one, and two vanish, while
\begin{equation}\label{jets:eq:q210}
 \sum_{u\in U_{210}}\Liord{3}{\zeta_{210}^u}
 =6\log2\log3\log5\log7.
\end{equation}

For the quadratic character $\chi_3$ with $\chi_3(1)=1$ and
$\chi_3(2)=-1$, the primitive trace at order one is
$\tau(\chi_3)L(1,\chi_3)=\iu\pi/3$. Since
$\chi_3(7)=\chi_3(13)=1$, one obtains
\begin{align}
 \sum_{u\in U_{21}}\chi_3(u)\Liord{1}{\zeta_{21}^u}
 &=-\frac{\iu\pi}{3}\log7,\label{jets:eq:chi3-21}\\
 \sum_{u\in U_{273}}\chi_3(u)\Liord{2}{\zeta_{273}^u}
 &=\frac{2\iu\pi}{3}\log7\log13.\label{jets:eq:chi3-273}
\end{align}
In the first equation the order-zero trace vanishes; in the second,
both order-zero and first-order traces vanish.
Likewise the character $\chi_4$ gives
\begin{equation}\label{jets:eq:chi4-20}
 \sum_{u\in U_{20}}\chi_4(u)\Liord{1}{\zeta_{20}^u}
 =-\frac{\iu\pi}{2}\log5.
\end{equation}
These primitive trace constants follow directly from
$\Li_1(z)=-\log(1-z)$ with the values continued from inside the unit
disk, so the signs do not depend on an unspecified logarithm branch.

For the real quadratic character $\chi_5$, positive on residues
$1,4$ and negative on $2,3$, let $\phi=(1+\sqrt5)/2$.
Pairing conjugates in the primitive trace gives
\[
 \cP_{5,\chi_5}(1)
 =2\log\frac{\sin(2\pi/5)}{\sin(\pi/5)}=2\log\phi.
\]
As $\chi_5(11)=1$,
\begin{equation}\label{jets:eq:chi5-55}
 \sum_{u\in U_{55}}\chi_5(u)\Liord{1}{\zeta_{55}^u}
 =-2\log\phi\log11.
\end{equation}
This derivation needs no numerical integer-relation search or
unproved independence statement about logarithms.

\subsection{A stable independent formula for numerical checks}
Differentiation of \eqref{jets:eq:polylog-Hurwitz-DFT} at its cancelled pole
gives, for $z=\zeta_q^u\ne1$,
\begin{equation}\label{jets:eq:stable-polylog-jet}
 \Liord{k}{z}
 =\frac{(-1)^k}{q}\sum_{a=1}^q z^a
 \sum_{j=0}^k\binom kj(\log q)^{k-j}\gamma_j(a/q).
\end{equation}
There is no pole term, because its complete Fourier coefficient
vanishes identically. This supplies a finite Stieltjes-coordinate
evaluator for the trace identities, independently of their
Euler-factor right sides. In development, direct tiny-step numerical
differentiation of a polylogarithm routine near integral order gave
an unreliable trace derivative; the pole-cancelled formula avoids
that cancellation problem. This observation is an implementation
diagnostic, not a claimed theorem about all versions of a library.

\subsection{A level-260 specialization and convention reconciliation}
We use the unbarred weight $\chi(a)$ in $\mathcal P_{q,\chi}$.
An alternative convention weights the trace by $\overline\chi(a)$; it
conjugates the character factors and interchanges the primitive $L$-character.
All formulas here use the first convention. The following exact specialization
adds two primes with character value one.



Let $\chi$ be primitive of conductor $f>1$, and $f\mid q$.
Among the added primes, define
\[
\mathcal S=\{p\mid q:p\nmid f,\ \chi(p)=1\},\qquad r_\chi=|\mathcal S|.
\]
The remaining added primes form $\mathcal N$.

\begin{theorem}[Twisted leading jet]\label{conductor:thm:twisted}
The trace $\cP_{q,\chi}$ vanishes to order at least $r_\chi$ at one, and
\begin{equation}\label{conductor:eq:twistedlead}
\cP_{q,\chi}^{(r_\chi)}(1)=
(-1)^{r_\chi}r_\chi!\,
\cP_{f,\chi}(1)
\prod_{p\in\mathcal S}\log p
\prod_{p\in\mathcal N}(1-\chi(p)).
\end{equation}
Its order is exactly $r_\chi$ whenever $\cP_{f,\chi}(1)\ne0$.
More generally its exact order is $r_\chi+\ord_{s=1}\cP_{f,\chi}(s)$.
\end{theorem}
\begin{proof}
The multiplier is
\[
B_{q,f,\chi}(1+\eps)=
(q/f)^{-\eps}
\prod_{\substack{p\mid q\\p\nmid f}}(1-\chi(p)p^\eps).
\]
The factors indexed by $\mathcal S$ have simple zeros with leading
coefficient $-\log p$. All other factors are nonzero at zero and have
value $1-\chi(p)$. Multiply their leading terms by the Taylor series
of $\cP_{f,\chi}(1+\eps)$.
\end{proof}

For primitive nonprincipal characters, the classical nonvanishing
$L(1,\overline\chi)\ne0$, together with the nonzero primitive Gauss sum in
\eqref{jets:eq:trace-formula}, supplies the stated nonvanishing hypothesis
\cite{jets:dlmf-dirichlet}. The explicit example below does not need that theorem;
its primitive trace is computed directly from logarithms.

For $\chi_4(a)=(-1)^{(a-1)/2}$ on odd integers,
\[
\cP_{4,\chi_4}(1)=\Li_1(\iu)-\Li_1(-\iu)=\frac{\iu\pi}{2}.
\]
Take $q=260=4\cdot5\cdot13$. Both added primes have $\chi_4(p)=1$.
Consequently
\begin{align}\label{conductor:eq:260}
\sum_{a\in\U_{260}}\chi_4(a)\Li_1(\zeta_{260}^a)&=0,\\
\sum_{a\in\U_{260}}\chi_4(a)\Liord{1}{\zeta_{260}^a}&=0,\\
\sum_{a\in\U_{260}}\chi_4(a)\Liord{2}{\zeta_{260}^a}
&=\iu\pi\log5\log13.
\end{align}
This is a proved identity, not an integer-relation conjecture. More generally,
for distinct primes $p_1,\ldots,p_r\equiv1\pmod4$ and
$q=4p_1\cdots p_r$,
\begin{equation}\label{conductor:eq:chi4family}
\sum_{a\in\U_q}\chi_4(a)\Liord{r}{\zeta_q^a}
=(-1)^r r!\frac{\iu\pi}{2}\prod_{j=1}^r\log p_j,
\end{equation}
with all lower order traces vanishing. The same leading identity persists
if the added prime factors occur with higher exponents.




\section{Explicit tables for the second constant}\label{tower:sec:tables}

We record the closed forms of $\gamma_2'=-[\zeta''(2,a)+2\zeta'(2,a)]$ (the
$s=2$ layer) and $\gamma_2''=2\zeta''(3,a)+6\zeta'(3,a)+2\zeta(3,a)$ (the $s=3$
layer) at the elementary denominators. The new constants entering are the
\emph{second} $s$-derivatives $\zeta''(2),\zeta''(3)$, $\beta''(s)=L''(s,\chim4)$
and $L''(s,\chim3)$; the values $L(3,\chim4)=\pi^3/32$ and
$L(3,\chim3)=4\pi^3/(81\sqrt3)$ are elementary. As at the first-derivative
layers, the odd-character block flips sign as a single unit under $p\mapsto q-p$
(including any elementary $\pi$-power value-part it carries). All entries are
exact and were verified to $<10^{-70}$.

\paragraph{Second constant, first parameter-derivative.}
\begin{align*}
\gamma_2'(1)   &= -\zeta''(2)-2\zeta'(2),\\
\gamma_2'(\tfrac12) &= -3\zeta''(2)-(6+8\ln2)\zeta'(2)-\tfrac43\pi^2\ln2-\tfrac23\pi^2\ln^2 2,\\
\gamma_2'(\tfrac13) &= -4\zeta''(2)-(8+9\ln3)\zeta'(2)-\tfrac32\pi^2\ln3-\tfrac34\pi^2\ln^2 3\\
 &\quad -\Bigl[\tfrac92 L''(2,\chim3)+9(1{+}\ln3)L'(2,\chim3)+\bigl(9\ln3+\tfrac92\ln^2 3\bigr)L(2,\chim3)\Bigr],\\
\gamma_2'(\tfrac14) &= -6\zeta''(2)-(12+28\ln2)\zeta'(2)-\tfrac{14}{3}\pi^2\ln2-5\pi^2\ln^2 2\\
 &\quad -\Bigl[8\beta''(2)+16(1{+}2\ln2)\beta'(2)+\bigl(32\ln2+32\ln^2 2\bigr)\Cat\Bigr],
\end{align*}
with $\gamma_2'(2/3),\gamma_2'(3/4)$ obtained by negating the bracketed
odd-character block, and the $q=6$ entries given in the source report.

\paragraph{Second constant, second parameter-derivative.}
\begin{align*}
\gamma_2''(1)   &= 2\zeta''(3)+6\zeta'(3)+2\zeta(3),\\
\gamma_2''(\tfrac12) &= 14\zeta''(3)+(42+32\ln2)\zeta'(3)+(14+48\ln2+16\ln^2 2)\zeta(3),\\
\gamma_2''(\tfrac13) &= 26\zeta''(3)+(78+54\ln3)\zeta'(3)+(26+81\ln3+27\ln^2 3)\zeta(3)\\
 &\quad +\Bigl[27L''(3,\chim3)+(81+54\ln3)L'(3,\chim3)+(27+81\ln3+27\ln^2 3)L(3,\chim3)\Bigr],\\
\gamma_2''(\tfrac14) &= 56\zeta''(3)+(168+240\ln2)\zeta'(3)+(56+360\ln2+248\ln^2 2)\zeta(3)\\
 &\quad +\Bigl[64\beta''(3)+(192+256\ln2)\beta'(3)+(2+12\ln2+8\ln^2 2)\pi^3\Bigr],
\end{align*}
again with $\gamma_2''(2/3),\gamma_2''(3/4)$ the odd-block negations. The
reducible fractions collapse correctly ($\tfrac24=\tfrac36=\tfrac12$,
$\tfrac26=\tfrac13$), an internal consistency check.

\section{The derivative\,$\leftrightarrow$\,negapolygamma duality}\label{tower:sec:duality}

The functional equation pairs the layer $s=k+1$ of $\gamma_1^{(k)}$ with the
layer $s=-k$, which is exactly where the negative-order polygamma
$\psi^{(-(k+1))}$ lives (its ladder uses $\zeta'(-k,\cdot)$). The seam is a
single uniform identity.

\begin{theorem}[Harmonic-number bridge law]\label{tower:thm:bridge}
Let $\chi$ be a primitive character mod $q$ for which $L(-k,\chi)\neq0$. Then
\begin{equation}\label{tower:eq:bridge}
  \frac{L'(k{+}1,\chi)}{L(k{+}1,\chi)}
   +\overline{\left(\frac{L'(-k,\chi)}{L(-k,\chi)}\right)}
   =\gamma+\ln\frac{2\pi}{q}-\h{k}.
\end{equation}
The same $\h{k}$ that sets the master coefficient in
Theorem~\ref{tower:thm:master} sets the bridge constant; the case $k=1$ recovers the
generalized-Glaisher constant $C_q=\gamma+\ln(2\pi/q)-1$.
\end{theorem}

\begin{proof}[Proof of the bridge]
For a primitive nonprincipal character of parity
$\chi(-1)=(-1)^\epsilon$, the completed function is
$\Lambda(s,\chi)=(q/\pi)^{(s+\epsilon)/2}
\Gamma((s+\epsilon)/2)L(s,\chi)$.
Its functional equation gives
$\Lambda'/\Lambda(s,\chi)=-\Lambda'/\Lambda(1-s,\bar\chi)$.
At $s=k+1$ with no trivial zero at $-k$, the two digamma terms,
using reflection and recurrence, sum to
$2(H_k-\gamma-\log2)$ before their factor $1/2$ is applied.
Collecting the conductor factor proves~\eqref{tower:eq:bridge}.
The conductor-one principal character requires the meromorphic completed
zeta function; the same nonzero-value calculation is valid for $k\ge1$
where neither point is a pole.
\end{proof}
\begin{remark}[Trivial-zero bridge]\label{tower:rem:parity}
For $k\ge1$, $L(-k,\chi)=0$ when $\epsilon\equiv k\pmod2$.
At a simple trivial zero, cancellation of the pole in the completed gamma
factor gives instead
\[
\frac{L'(k+1,\chi)}{L(k+1,\chi)}+
\overline{\frac{L''(-k,\chi)}{2L'(-k,\chi)}}
=\gamma+\log(2\pi/q)-H_k.
\]
The factor $2$ in the denominator is essential. This follows by writing
$L(-k+u)=L'(-k)u+\frac12L''(-k)u^2+O(u^3)$ in the same
logarithmic functional equation and taking its finite part.
\end{remark}

At $z=\tfrac14$ the governing character is $\chim4$ (odd). For $k$ even there is
no trivial zero, and $\psi^{(-(k+1))}(\tfrac14)$ carries the unreduced derivative coordinate $\beta'(k{+}1)=L'(k{+}1,\chim4)$, so $\gamma_1^{(k)}(\tfrac14)$ is an exact
finite combination of $\psi^{(-(k+1))}(\tfrac14)$ and the remaining zeta-value and zeta-derivative terms. The
representative instance $k=2$ is, after eliminating $\beta'(3)$ entirely,
\begin{equation}\label{tower:eq:k2duality}
\begin{aligned}
  \gamma_1''(\tfrac14)=\;&-256\,\pi^3\,\psi^{(-3)}(\tfrac14)
   -56\,\zeta'(3)-84\,\zeta(3)-120\ln2\,\zeta(3)+35\pi\,\zeta(3)\\
   &-3\pi^3-4\ln2\,\pi^3+64\ln\Glaisher\,\pi^3-2\gamma\,\pi^3+2\ln(2\pi)\,\pi^3 ,
\end{aligned}
\end{equation}
verified to $<10^{-118}$. (The irrational coefficient $-256\pi^3$ on
$\psi^{(-3)}(\tfrac14)$ is why a rational integer-relation search for this
identity must take $\pi^3\psi^{(-3)}$, not $\psi^{(-3)}$, as the search atom.)
For $k$ odd the value-level duality \emph{degenerates}: by
Remark~\ref{tower:rem:parity}, $\psi^{(-(k+1))}(\tfrac14)$ then carries only the lower
\emph{value} atom (for $k=1$, $\beta'(-1)=2\Cat/\pi$, a value rather than a derivative at the positive point; for $k=3$,
$\beta'(-3)\sim\beta(4)/\pi^3$), while the genuine derivative atom
$\beta'(k{+}1)$ of $\gamma_1^{(k)}(\tfrac14)$ is tied instead to the second
derivative $\beta''(-k)$, absent from $\psi^{(-(k+1))}$. Numerically the
coefficient of $\beta'(k{+}1)$ in $\psi^{(-(k+1))}(\tfrac14)$ is zero to $\sim
10^{-163}$ for $k=1,3$. The clean odd-$k$ value-level duality therefore lives at
even-character points (e.g.\ $q=5$), where the parity is reversed.

\section{One normalized functional equation for all derivative bridges}
\label{jets:sec:bridges}

The Stieltjes drafts connect derivatives of \(L(s,\chi)\) at a positive
integer with derivatives at its reflected negative integer. A single
identity of analytic germs organizes all these formulas and treats
trivial zeros without differentiating singular logarithms. It also
makes the conjugation required for complex characters unavoidable.
The input is the classical primitive Dirichlet functional equation
and the reflection and duplication formulas for \(\Gamma\)
\cite{hstruct:DLMF}; the all-order normalized form below is a systematic
continuation of the bridges used in the corpus.

\subsection{Regularization and the normalized germ}

Let \(\chi\) be primitive of conductor \(q\), allowing the trivial
character of conductor \(1\), for which \(L(s,\chi)=\zeta(s)\).
Fix \(k\geq1\), and write
\[
 \delta=\frac{1-\chi(-1)}2\in\{0,1\},\qquad
 \nu=
 \begin{cases}
 1,&\chi(-1)=(-1)^k,\\
 0,&\chi(-1)\ne(-1)^k.
 \end{cases}
\]
Define
\begin{equation}
 B(t)=\frac{L(-k+t,\chi)}{t^\nu},
 \qquad
 B^\#(z)=\overline{B(\overline z)},
 \label{jets:eq:bridge-regularization}
\end{equation}
with the removable value filled in when \(\nu=1\).
The function \(B^\#\) is holomorphic near zero; it is coefficientwise
conjugation of the Taylor series of \(B\).

\begin{theorem}[Normalized positive--negative germ]
\label{jets:thm:bridge-normalized}
The quantities \(L(k+1,\chi)\) and \(B(0)\) are nonzero. On a
sufficiently small disk about \(z=0\),
\begin{equation}
 \frac{L(k+1+z,\chi)}{L(k+1,\chi)}
 =
 K_\nu(z)\frac{B^\#(-z)}{B^\#(0)},
 \label{jets:eq:bridge-normalized-germ}
\end{equation}
where
\begin{equation}
 \begin{aligned}
 K_\nu(z)&=
 \left(\frac{2\pi}{q}\right)^z
 \frac{\Gamma(k+1)}{\Gamma(k+1+z)}\,T_\nu(z),\\
 T_0(z)&=\sec\frac{\pi z}{2},\qquad
 T_1(z)=\frac{\pi z/2}{\sin(\pi z/2)},\quad T_1(0)=1.
 \end{aligned}
 \label{jets:eq:bridge-elementary-kernel}
\end{equation}
All three normalized factors in
\eqref{jets:eq:bridge-normalized-germ} equal \(1\) at the origin.
\end{theorem}

\begin{proof}
The Euler product gives \(L(k+1,\chi)\ne0\), since \(k+1>1\).
Use the completed function
\[
 \Lambda(s,\chi)=
 \left(\frac q\pi\right)^{(s+\delta)/2}
 \Gamma\!\left(\frac{s+\delta}{2}\right)L(s,\chi)
\]
and its functional equation
\begin{equation}
 \Lambda(s,\chi)=\varepsilon_\chi\Lambda(1-s,\overline\chi),
 \qquad |\varepsilon_\chi|=1.
 \label{jets:eq:bridge-completed-fe}
\end{equation}
At \(s=-k\), the gamma factor is finite and nonzero if
\(\nu=0\), and has a simple pole if \(\nu=1\).
The reflected completed value at \(k+1\) is finite and nonzero.
Consequently \(L(-k,\chi)\ne0\) in the first case, and
\(L(s,\chi)\) has a simple zero at \(-k\) in the second.
In either case \(B\) is holomorphic and \(B(0)\ne0\).

Put
\[
 a=\frac{\delta-k}{2},\qquad
 b=\frac{k+1+\delta}{2}.
\]
Solving \eqref{jets:eq:bridge-completed-fe} at \(s=k+1+z\) gives
\begin{equation}
 L(k+1+z,\chi)
 =\varepsilon_\chi
 \left(\frac q\pi\right)^{-k-\frac12-z}
 \frac{\Gamma(a-z/2)}{\Gamma(b+z/2)}
 L(-k-z,\overline\chi).
 \label{jets:eq:bridge-fe-before-normalization}
\end{equation}
Since
\(L(-k-z,\overline\chi)=(-z)^\nu B^\#(-z)\),
the numerator
\[
 F_\nu(z)=(-z)^\nu\Gamma(a-z/2)
\]
has a nonzero removable value at the origin.
Dividing \eqref{jets:eq:bridge-fe-before-normalization} by its value
at \(z=0\) cancels the root number and all constant factors:
\begin{equation}
 \frac{L(k+1+z,\chi)}{L(k+1,\chi)}
 =
 \left(\frac q\pi\right)^{-z}
 \frac{F_\nu(z)}{F_\nu(0)}
 \frac{\Gamma(b)}{\Gamma(b+z/2)}
 \frac{B^\#(-z)}{B^\#(0)}.
 \label{jets:eq:bridge-normalized-intermediate}
\end{equation}

If \(\nu=0\), then \(a\) is a half-integer.
The gamma reflection formula and
\(\sin(\pi(a-z/2))=\sin(\pi a)\cos(\pi z/2)\)
give
\begin{equation}
 \frac{F_0(z)}{F_0(0)}
 =\frac{\Gamma(1-a)}{\Gamma(1-a+z/2)}
       \sec\frac{\pi z}{2}.
 \label{jets:eq:bridge-gamma-half-integer}
\end{equation}
If \(\nu=1\), write \(a=-m\), \(m\geq0\).
Then \(F_1(0)=2(-1)^m/m!\), and the same reflection formula
gives
\begin{equation}
 \frac{F_1(z)}{F_1(0)}
 =\frac{\Gamma(1-a)}{\Gamma(1-a+z/2)}
       \frac{\pi z/2}{\sin(\pi z/2)}.
 \label{jets:eq:bridge-gamma-integer}
\end{equation}
For example, the pole in \(\Gamma(-m-z/2)\) is canceled by
the explicit factor \(-z\), so no logarithm of a function
vanishing at zero has been taken.

The unordered pair \(\{b,1-a\}\) is
\(\{(k+1)/2,(k+2)/2\}\), for either value of \(\delta\).
Duplication therefore yields
\begin{equation}
 \frac{\Gamma(b)\Gamma(1-a)}
 {\Gamma(b+z/2)\Gamma(1-a+z/2)}
 =2^z\frac{\Gamma(k+1)}{\Gamma(k+1+z)}.
 \label{jets:eq:bridge-gamma-duplication}
\end{equation}
Substituting \eqref{jets:eq:bridge-gamma-half-integer} or
\eqref{jets:eq:bridge-gamma-integer}, followed by
\eqref{jets:eq:bridge-gamma-duplication}, into
\eqref{jets:eq:bridge-normalized-intermediate} proves the theorem.
\end{proof}

\begin{remark}
The statement concerns analytic germs and local logarithms.
It requires no claim about a global zero-free region.
For a complex character the functional equation involves
\(\overline\chi\), and therefore the reflected Taylor coefficients
are conjugated. Replacing \(B^\#\) by \(B\) is justified for a real
character but generally false for a complex one.
\end{remark}

\subsection{Closed logarithmic coefficients at every order}

Choose the normalized local logarithms, each zero at the origin,
of \(L(k+1+z,\chi)/L(k+1,\chi)\) and \(B(t)/B(0)\).
For \(r\geq1\), set
\[
 \mathcal L_r=
 \left.\frac{d^r}{ds^r}\log L(s,\chi)\right|_{s=k+1},
 \qquad
 b_r=
 \left.\frac{d^r}{dt^r}\log B(t)\right|_{t=0}.
\]
These derivatives do not depend on a choice of logarithm branch.
Write
\[
 H_k^{(r)}=\sum_{j=1}^{k}j^{-r},\qquad H_k=H_k^{(1)},
 \qquad
 \eta(r)=(1-2^{1-r})\zeta(r)\quad(r\geq2).
\]

\begin{theorem}[All-order logarithmic bridges]
\label{jets:thm:bridge-log-jets}
For every \(r\geq1\),
\begin{equation}
 \mathcal L_r-(-1)^r\overline{b_r}=\kappa_r,
 \label{jets:eq:bridge-log-jets}
\end{equation}
where
\begin{equation}
 \kappa_r=
 \begin{cases}
 \displaystyle\gamma+\log\frac{2\pi}{q}-H_k,
       &r=1,\\[4pt]
 \displaystyle(r-1)!\bigl(\zeta(r)-H_k^{(r)}\bigr),
       &r\geq3\ \text{odd},\\[4pt]
 \displaystyle(r-1)!\bigl(H_k^{(r)}+(-1)^\nu\eta(r)\bigr),
       &r\geq2\ \text{even}.
 \end{cases}
 \label{jets:eq:bridge-kappa}
\end{equation}
\end{theorem}

\begin{proof}
Taking normalized logarithms of
\eqref{jets:eq:bridge-normalized-germ} gives
\[
 \log\frac{L(k+1+z,\chi)}{L(k+1,\chi)}
 =\log K_\nu(z)
  +\overline{\log\frac{B(-\overline z)}{B(0)}}.
\]
The \(r\)-th derivative of the last term at zero is
\((-1)^r\overline{b_r}\). It remains to compute the
derivatives of \(\log K_\nu\).

The convergent gamma expansion, on \(|z|<1\), is
\begin{equation}
 \log\frac{\Gamma(k+1+z)}{\Gamma(k+1)}
 =(H_k-\gamma)z+
 \sum_{r\geq2}\frac{(-1)^r}{r}
       \bigl(\zeta(r)-H_k^{(r)}\bigr)z^r.
 \label{jets:eq:bridge-log-gamma-series}
\end{equation}
The sine and cosine products give, in the same disk,
\begin{align}
 \log\sec\frac{\pi z}{2}
 &=\sum_{m\geq1}
   \frac{(1-2^{-2m})\zeta(2m)}{m}z^{2m},
 \label{jets:eq:bridge-log-sec-series}\\
 \log\frac{\pi z/2}{\sin(\pi z/2)}
 &=\sum_{m\geq1}
   \frac{2^{-2m}\zeta(2m)}{m}z^{2m}.
 \label{jets:eq:bridge-log-sinc-series}
\end{align}
These identities also follow directly by taking logarithms
of the absolutely locally convergent products
\(\cos(\pi z/2)=\prod_{j\geq0}(1-z^2/(2j+1)^2)\) and
\(\sin(\pi z/2)/(\pi z/2)
=\prod_{j\geq1}(1-z^2/(2j)^2)\).

The linear term of \(\log K_\nu\) is the stated
\(\kappa_1\). Odd coefficients of order at least three come
only from the negative of
\eqref{jets:eq:bridge-log-gamma-series}, giving the second case.
For \(r=2m\), the secant coefficient combines with the gamma
coefficient to give
\[
 (r-1)!\left[
  H_k^{(r)}-\zeta(r)+2(1-2^{-r})\zeta(r)\right]
 =(r-1)!\bigl(H_k^{(r)}+\eta(r)\bigr).
\]
Replacing the secant factor by the sine quotient instead gives
\[
 (r-1)!\left[
  H_k^{(r)}-\zeta(r)+2^{1-r}\zeta(r)\right]
 =(r-1)!\bigl(H_k^{(r)}-\eta(r)\bigr).
\]
These are exactly the two even cases in
\eqref{jets:eq:bridge-kappa}.
\end{proof}

At a nonzero reflected value, \(\nu=0\), the first two identities
read
\begin{align}
 \frac{L'(k+1,\chi)}{L(k+1,\chi)}
 +\overline{\frac{L'(-k,\chi)}{L(-k,\chi)}}
 &=\gamma+\log\frac{2\pi}{q}-H_k,
 \label{jets:eq:bridge-first-nonzero}\\
 (\log L)''(k+1,\chi)
 -\overline{(\log L)''(-k,\chi)}
 &=H_k^{(2)}+\frac{\pi^2}{12}.
 \label{jets:eq:bridge-second-nonzero}
\end{align}
At a trivial zero, \(\nu=1\), the Taylor expansion
\[
 B(t)=L'(-k,\chi)
   +\frac{L''(-k,\chi)}2t
   +\frac{L'''(-k,\chi)}6t^2+\cdots
\]
instead gives
\begin{align}
 \frac{L'(k+1,\chi)}{L(k+1,\chi)}
 +\frac12\overline{\frac{L''(-k,\chi)}{L'(-k,\chi)}}
 &=\gamma+\log\frac{2\pi}{q}-H_k,
 \label{jets:eq:bridge-first-zero}\\
 (\log L)''(k+1,\chi)
 -\overline{\left(
    \frac{L'''(-k,\chi)}{3L'(-k,\chi)}
   -\frac{L''(-k,\chi)^2}{4L'(-k,\chi)^2}
   \right)}
 &=H_k^{(2)}-\frac{\pi^2}{12}.
 \label{jets:eq:bridge-second-zero}
\end{align}
These formulas preserve the familiar first-order harmonic
bridge and supply a uniform extension to every derivative
order and both character parities.

\subsection{Raw jets and an exact triangular implementation}

The normalized germ also gives direct formulas for ordinary
derivatives, without first converting each order to logarithmic
derivatives. Define \(K_m\) by
\[
 K_\nu(z)=\sum_{m\geq0}K_m\frac{z^m}{m!},\qquad K_0=1.
\]
Since \(\log K_\nu(z)=\sum_{r\geq1}\kappa_rz^r/r!\),
these coefficients are the complete exponential Bell polynomials:
\begin{equation}
 K_m=
 m!\!\!\sum_{\substack{m_1,\ldots,m_m\geq0\\
                     m_1+2m_2+\cdots+mm_m=m}}
 \prod_{r=1}^{m}\frac1{m_r!}
       \left(\frac{\kappa_r}{r!}\right)^{m_r}
 \qquad(m\geq1).
 \label{jets:eq:bridge-bell-coefficients}
\end{equation}
Leibniz's rule in \eqref{jets:eq:bridge-normalized-germ} proves
\begin{equation}
 \frac{L^{(m)}(k+1,\chi)}{L(k+1,\chi)}
 =
 \sum_{j=0}^{m}\binom mj K_{m-j}(-1)^j
       \overline{\frac{B^{(j)}(0)}{B(0)}}.
 \label{jets:eq:bridge-raw-jets}
\end{equation}
The required regularized derivatives are explicitly
\begin{equation}
 \frac{B^{(j)}(0)}{B(0)}
 =\frac{j!\,\nu!}{(j+\nu)!}
    \frac{L^{(j+\nu)}(-k,\chi)}{L^{(\nu)}(-k,\chi)}.
 \label{jets:eq:bridge-regularized-raw-jets}
\end{equation}
The diagonal coefficient in
\eqref{jets:eq:bridge-raw-jets} is \((-1)^m\), so the transformation
is triangular and invertible. Its coefficients belong to the
ring generated by the explicit \(\kappa_r\)'s. This supplies
exact relation certificates at arbitrary order and avoids
repeated symbolic differentiation across gamma poles.

\begin{remark}[Correction of the complex-character formulas]
The older derivative-relations report omits the conjugation
in versions of the trivial-zero first bridge and the nonzero
second bridge. Equations~\eqref{jets:eq:bridge-first-zero} and
\eqref{jets:eq:bridge-second-nonzero} give the required replacements.
The two main Stieltjes articles already include conjugation
in their displayed first bridges; those correct formulas
should be retained. Exact source locators and numerical
counterexamples to the omitted-conjugation versions appear
in the correction register.
\end{remark}


\section{Experimental independence tests}
The tables introduce $\zeta''(2),\zeta''(3),\beta''(2),\beta''(3)$ and
$L''(2,\chi_{-3}),L''(3,\chi_{-3})$. Historical searches at $200$ digits
with planted controls recovered the known value reductions and found no
accepted relation for the six derivative coordinates in the lower-order
baskets. This supports using them as named coordinates in a reduction
table. Their mutual independence, and the impossibility of reductions in
larger baskets, remain open.

\section{Elementary grids and the sixth-point completion}
For $q=3,4,6$ define
\[
A_3(s)=\frac{3^s-1}{2}\zeta(s),\quad
A_4(s)=\frac{4^s-2^s}{2}\zeta(s),\quad
A_6(s)=\frac{(2^s-1)(3^s-1)}2\zeta(s),
\]
\[
B_3(s)=\frac{3^s}{2}L(s,\chi_{-3}),\quad
B_4(s)=\frac{4^s}{2}\beta(s),\quad
B_6(s)=\frac{6^s+3^s}{2}L(s,\chi_{-3}).
\]
Residue-class decomposition proves
$\zeta(s,1/q)=A_q(s)+B_q(s)$ and
$\zeta(s,1-1/q)=A_q(s)-B_q(s)$ for these three $q$.
Consequently all four derivative tables, including the omitted sixth-point
rows in the later article, follow from the differential operators
\begin{align*}
\gamma_1'(a)&=(\partial_s+1)\zeta(s,a)|_{s=2},&
\gamma_1''(a)&=-(2\partial_s+3)\zeta(s,a)|_{s=3},\\
\gamma_2'(a)&=-(\partial_s^2+2\partial_s)\zeta(s,a)|_{s=2},&
\gamma_2''(a)&=(2\partial_s^2+6\partial_s+2)\zeta(s,a)|_{s=3}.
\end{align*}
This supplies an exact symbolic recipe for every entry, without new fitted
constants. Write $a=\log2$, $b=\log3$, and $L_j(s)=L^{(j)}(s,\chi_{-3})$.
At the sixth points the first-constant rows are
\begin{align*}
\gamma_1'(1/6),\ \gamma_1'(5/6)
&=12\zeta'(2)+2\pi^2+\frac83\pi^2a+\frac94\pi^2b\\
&\quad\pm\left[\frac{45}{2}L_1(2)
+\left(\frac{45}{2}+18a+\frac{45}{2}b\right)L_0(2)\right],\\
\gamma_1''(1/6),\ \gamma_1''(5/6)
&=-182\zeta'(3)-(273+208a+189b)\zeta(3)\\
&\quad\mp\left[243L_1(3)
+\left(\frac{729}{2}+216a+243b\right)L_0(3)\right].
\end{align*}
The second-constant rows are
\begin{align*}
\gamma_2'(1/6),\ \gamma_2'(5/6)&=S_2\mp O_2,\\
S_2&=-12\zeta''(2)-(24+32a+27b)\zeta'(2)\\
&\quad-\pi^2\left(\frac{16}3a+\frac92b+
\frac83a^2+\frac94b^2+6ab\right),\\
O_2&=\frac{45}{2}L_2(2)+(45+36a+45b)L_1(2)\\
&\quad+\left(36a+45b+18a^2+\frac{45}{2}b^2+36ab\right)L_0(2),\\
\gamma_2''(1/6),\ \gamma_2''(5/6)&=S_3\pm O_3,\\
S_3&=182\zeta''(3)+(546+416a+378b)\zeta'(3)\\
&\quad+(182+624a+567b+208a^2+189b^2+432ab)\zeta(3),\\
O_3&=243L_2(3)+(729+432a+486b)L_1(3)\\
&\quad+(243+648a+729b+216a^2+243b^2+432ab)L_0(3).
\end{align*}
The upper sign is the $1/6$ entry. These formulas consolidate the tables
from the first-derivative, second-derivative and general-order reports.


% Collective reconciliation of two independent periodic-contact routes.
\providecommand{\pccD}{\mathrm D}
\providecommand{\pccZ}{\mathcal Z}
\providecommand{\pccA}{\mathcal A}
\providecommand{\pccE}{\mathcal E}
\providecommand{\pccJ}{\mathsf J}
\providecommand{\pccC}{\mathsf C}
\providecommand{\pccI}{\mathcal I}
\section{Fixed-coordinate periodic collision contacts}
\label{pc:scope}
The differentiation anomaly of Theorem~\ref{stconv:thm:anomaly} already
fixes the individual periodic factors. A second question remains: does
convolution of those factors equal the coordinate finite-part extension
of their separated correlation? The independent continuations
\cite{pcc:reports} give the same answer at every index and derivative order.
The comparison below keeps the coordinate $x$ for each original factor and
the independent coordinates $a$ and $1-a$ for the shift. In this convention
there is exactly one contact term. Its coefficient is distinct from the
scalar constant in an off-point collision expansion.

Write $T_{m,p}=F_{m,p}$, with $T_{m,0}=G_m$, in the notation of
\eqref{stconv:eq:Fdefinition}. Reflection is denoted by a check, so
$\widehat{\check U}(k)=\widehat U(-k)$. Set
\[
 P_p(u)=\frac{(1+u)_p}{p!}=\prod_{j=1}^p(1+u/j),\qquad P_0=1,
 \qquad H_p^{(j)}=\sum_{l=1}^p l^{-j}.
\]
All indices $m,n,p,q$ are nonnegative integers; put $r=p+q$ and $w=u+v$.
The scalar Stieltjes germ and its argument derivatives are
\begin{align}
 g(u,x)&=\zeta(1+u,x)-1/u
       =\sum_{m\ge0}(-1)^m\gamma_m(x)u^m/m!,\label{pc:gdef}\\
 G_p(u,x)&=\partial_x^p g(u,x)
   =(-1)^p(1+u)_p\zeta(1+p+u,x)-\delta_{p0}/u.\label{pc:G}
\end{align}
The apparent pole in $g$ is removable for each $x>0$.

\subsection{Whole resonant numerators and Fourier input}
Let $\pccZ_s=Z(s)$ be the meromorphic family of
Lemma~\ref{stconv:lem:Laurentdist}. Its Fourier coefficients are
\begin{equation}\label{pc:hurwitzFourier}
 \widehat{\pccZ_s}(0)=0,\qquad
 \widehat{\pccZ_s}(k)=\Gamma(1-s)(2\pi\iu k)^{s-1}\quad(k\ne0),
\end{equation}
on the branch \eqref{stconv:eq:logbranch}.
\begin{lemma}[Full resonant completion]\label{pc:resonant}
The germ generating the unit-coordinate finite parts is
\begin{equation}\label{pc:Tpmeromorphic}
 T_p(u):=\sum_{m\ge0}\frac{(-1)^m}{m!}T_{m,p}u^m
 =\pccD^p\pccZ_{1+u}-\frac{\delta_{p0}}u\,1
       +\frac{P_p(u)}u\delta_0^{(p)}.
\end{equation}
It is holomorphic near zero, and every $T_{m,p}$ has mean zero.
\end{lemma}
\begin{proof}
The local singular part of $G_p$ is
$c_p(u)x^{-p-1-u}$, with $c_p(u)=(-1)^pp!P_p(u)$; its remaining
shifted Hurwitz part is jointly holomorphic. Subtract the test-function
Taylor polynomial through degree $p$. Its indices $j<p$ contribute
$c_p(u)\varphi^{(j)}(0)/(j!(j-p-u))$, holomorphic near zero.
The index $p$ contributes
\[
 -\frac{c_p(u)\varphi^{(p)}(0)}{p!u}
   =-\frac{P_p(u)}u\langle\delta_0^{(p)},\varphi\rangle.
\]
Every coefficient of this resonant monomial is proportional to
$x^{-1}\log^j x$, whose unit-coordinate finite-part integral on $(0,1)$
is zero. Thus its whole numerator must be removed, including its regular
Taylor coefficients. This gives \eqref{pc:Tpmeromorphic}, using
$\pccD^p\pccZ_{1+u}=(-1)^p(1+u)_p\pccZ_{1+p+u}$.
After subtraction the local remainder, with any fixed number of spectral
derivatives, is dominated by $C x^{-\eta}(1+|\log x|^M)$, $\eta<1$.
Consequently pairing and coefficient extraction are holomorphic and give
exactly the cutoff finite parts. The nonlocal zero Fourier mode is zero;
for $p=0$ the last two terms cancel in mean, and for $p>0$ both have
mean zero.
\end{proof}

The existing differentiation anomaly becomes
\begin{equation}\label{pc:derivativeformula}
 T_{m,p}=\pccD^pG_m+\kappa_{m,p}\delta_0^{(p)},\qquad
 \kappa_{m,p}=(-1)^m m![u^{m+1}]P_p(u).
\end{equation}
This is the earlier elementary-symmetric formula, since
$[u^{m+1}]P_p=e_{m+1}(p)$. It also follows immediately by subtracting
$\pccD^pT_0(u)$ from \eqref{pc:Tpmeromorphic}. In particular
$\kappa_{0,p}=H_p$ and $\kappa_{m,p}=0$ when $m\ge p$.

For a scalar complete exponential Bell polynomial $B_j$, define
\begin{equation}\label{pc:bell}
 b_m(z)=\frac{(-1)^{m+1}}{m+1}
 B_{m+1}(z,1!\zeta(2),2!\zeta(3),\ldots,m!\zeta(m+1)).
\end{equation}
Equivalently, $B_j$ is defined by
$\exp(\sum_{l\ge1}x_lu^l/l!)=\sum_{j\ge0}B_ju^j/j!$.
For $k\ne0$, put $L_k=\gamma+\log(2\pi\iu k)$. Then
\begin{equation}\label{pc:Tfourier}
 \widehat T_{m,p}(k)=(2\pi\iu k)^p
             \bigl(b_m(L_k)+\kappa_{m,p}\bigr),\qquad
 \widehat T_{m,p}(0)=0.
\end{equation}
Indeed \eqref{pc:Tpmeromorphic} gives the nonzero-mode germ
$(2\pi\iu k)^p\{\Gamma(-u)(2\pi\iu k)^u+P_p(u)/u\}$.
Insert $\Gamma(-u)=-\Gamma(1-u)/u$ and
\[
 \Gamma(1-u)(2\pi\iu k)^u
 =\exp\left(L_ku+\sum_{j\ge2}\zeta(j)u^j/j\right),
\]
the classical gamma expansion \cite{pcc:DLMFGamma}. This proves
\eqref{pc:Tfourier}; for example $b_0(z)=-z$ and
$b_1(z)=(z^2+\zeta(2))/2$.

The lowest-index comparison, proved in the general theorem below, is
\begin{equation}\label{pc:preview}
 \check G_0*G_0
 =\FP_{\stconvT,a}\left[\gamma_1(a)+\gamma_1(1-a)-\pi^2/3\right]
       +\frac{\pi^2}{3}\delta_0.
\end{equation}
Both factors have zero mean. The atom restores that same mean in the
convolution; it is invisible in the punctured-circle function.
\section{The off-collision kernel and its second finite part}\label{pc:offpoint}
For $0<a<1$, set $b=1-a$ and define, in the original integration variable,
\begin{equation}\label{pc:Jpoint}
 J_{mn}^{p,q}(a)=\FP\int_0^1
 \gamma_m^{(p)}(x)\gamma_n^{(q)}(\{x+a\})\dd x.
\end{equation}
There are two independent right-hand singularities, at $x=0+$ and $x=b+$; the second coordinate is $x-b$. This is the convention of \cite{pcc:reports}. Define the distinct distribution
\begin{equation}\label{pc:Jdistribution}
 \pccJ_{mn}^{p,q}=\FP_{\stconvT,a}J_{mn}^{p,q}(a).
\end{equation}
The second finite part is taken in the shift variable, at both $a=0+$ and $a=1-$.

We recall the off-point generating kernel, deriving the ingredients needed for the distributional comparison. Put
\begin{align}
 \pccE(u,v)&=\frac{\Gamma(1-u)\Gamma(1+u+v)}{\Gamma(1+v)},\label{pc:E}\\
 \pccA(u,v)&=\frac{\Gamma(1-u)\Gamma(1-v)}{\Gamma(1-u-v)}
 \frac{\cos(\pi(u-v)/2)}{\cos(\pi(u+v)/2)}.\label{pc:A}
\end{align}
These are holomorphic near the origin and satisfy
\begin{equation}\label{pc:EA}
 v\pccE(u,v)+u\pccE(v,u)=(u+v)\pccA(u,v),\qquad
 \pccA(u,0)=\pccA(0,v)=1.
\end{equation}
The identity follows by inserting the gamma reflection formula. Equivalently, it follows by comparing the two Hurwitz forms of the classical spectral correlation.

For $r\ge1$, write
\[
 W_r(t,x)=(1+t)_r\zeta(1+r+t,x)=(-1)^rG_r(t,x).
\]
The off-point germ whose coefficients are $(-1)^{m+n}J_{mn}^{p,q}/(m!n!)$ is
\begin{align}\label{pc:Jgerm}
 \mathcal J_{pq}(u,v;a)
 ={}&(-1)^q\frac{P_p(u)W_r(v,a)-\pccE(u,v)W_r(w,a)}u\\
 &+(-1)^p\frac{P_q(v)W_r(u,b)-\pccE(v,u)W_r(w,b)}v.\notag
\end{align}
For $p=q=0$, the corresponding expression is
\begin{align}\label{pc:Jgerm0}
 \mathcal J_{00}(u,v;a)
 ={}&\frac{g(v,a)-\pccE(u,v)g(w,a)}u
   +\frac{g(u,b)-\pccE(v,u)g(w,b)}v\notag\\
 &+\frac{1-\pccA(u,v)}{uv}.
\end{align}
Each divided difference is holomorphic: its numerator vanishes on the relevant spectral axis.

To verify these formulas, start where all ordinary integrals converge and use the Hurwitz Fourier series. The shifted spectral kernel is
\begin{align}\label{pc:beta}
 C(s,t;a)={}&B(1-s,s+t-1)\zeta(s+t-1,a)\notag\\
 &+B(1-t,s+t-1)\zeta(s+t-1,1-a).
\end{align}
It follows either by gamma reflection from the product of Fourier coefficients, or from the classical polylogarithmic expression. This product-kernel method is classical \cite{stieltjes:EspinosaMoll1}; the shifted specialization and its Stieltjes completion are already used in \cite{pcc:reports}. Multiply by $(-1)^{p+q}(1+u)_p(1+v)_q$, set $s=1+p+u,t=1+q+v$, and remove the \emph{full} resonant Taylor contribution at each of the two original singularities. The gamma identity
\[
 (-1)^p(1+u)_p\Gamma(-p-u)=\Gamma(-u)
\]
then gives \eqref{pc:Jgerm}; the subtractions of $1/u$ and $1/v$ give \eqref{pc:Jgerm0}. The same finite Taylor-subtraction argument as in Lemma~\ref{pc:resonant} justifies holomorphy and coefficient extraction.

The exact Hurwitz shift shows that, near $a=0+$,
\[
 J_{mn}^{p,q}(a)=a^{-r-1}\mathcal Q_{mn}^{p,q}(\log a)
                         +\text{an analytic function of }a,
\]
where $\deg\mathcal Q_{mn}^{p,q}\le m+n+1$. The other endpoint is obtained by interchanging $(m,p)$ and $(n,q)$ and reflecting. Consequently \eqref{pc:Jdistribution} is well defined, with only finitely many test-function Taylor subtractions at each endpoint.

Taking the shift-variable finite part in \eqref{pc:Jgerm} yields, for $r\ge1$,
\begin{align}\label{pc:Jfpgen}
 \pccJ_{pq}(u,v)
 ={}&(-1)^p\frac{P_p(u)T_r(v)-\pccE(u,v)T_r(w)}u\notag\\
 &+(-1)^q\frac{P_q(v)\check T_r(u)-\pccE(v,u)\check T_r(w)}v.
\end{align}
For $r=0$, use the same formula with $P_0=1$ and add $(1-\pccA)/(uv)$ times the constant distribution $1$. This expression specifies the full second finite part; it is not merely a choice of extension having the correct off-point values.

\section{The all-index collision-contact theorem}\label{pc:mainsection}
Define the convolution distribution
\begin{equation}\label{pc:Cdef}
 \pccC_{mn}^{p,q}=\check T_{m,p}*T_{n,q}.
\end{equation}
Its restriction to the punctured circle is \eqref{pc:Jpoint}. Indeed, for a nonzero shift the singular supports of the two factors are disjoint. Multiplication by the other, smooth factor near each singular point gives exactly the two original coordinate finite parts. A partition of unity in $x$ makes this assertion a direct consequence of the distribution definitions. It does not involve defining a pointwise product of two distributions at a coincident singularity.

\begin{theorem}[Complete periodic contact correction]\label{pc:main}
For all nonnegative $m,n,p,q$, with $r=p+q$,
\begin{equation}\label{pc:mainformula}
 \boxed{\ \pccC_{mn}^{p,q}=\pccJ_{mn}^{p,q}
                     +\Delta_{mn}^{p,q}\delta_0^{(r)}.\ }
\end{equation}
The coefficient is
\begin{align}\label{pc:deltacoeff}
 \Delta_{mn}^{p,q}={}&(-1)^{m+n+p}m!n![u^{m+1}v^{n+1}]\mathcal N_{pq}(u,v),\\
 \mathcal N_{pq}(u,v)={}&\pccA(u,v)P_r(u+v)-P_p(u)P_r(v)\notag\\
 &-P_q(v)P_r(u)+P_p(u)P_q(v).\label{pc:N}
\end{align}
Equivalently, the spectral generating coefficient is the holomorphic germ
\begin{equation}\label{pc:Deltagen}
 \Delta_{pq}(u,v)=\frac{(-1)^p}{uv}\mathcal N_{pq}(u,v).
\end{equation}
There are no additional lower-order delta derivatives in this fixed convention.
\end{theorem}
\begin{proof}
First let $r\ge1$. Work with meromorphic distribution germs before taking finite parts. Set $U_p(u)=\pccD^p\pccZ_{1+u}$. By \eqref{pc:beta} and the gamma cancellation preceding \eqref{pc:Jfpgen}, the raw correlation is
\[
 \check U_p(u)*U_q(v)
 =-\frac{(-1)^p\pccE(u,v)}u\pccD^r\pccZ_{1+w}
  -\frac{(-1)^q\pccE(v,u)}v\check{\pccD^r\pccZ_{1+w}}.
\]
This identity can also be checked at each nonzero Fourier mode from \eqref{pc:hurwitzFourier}. It holds as a meromorphic distribution identity because it holds in an ordinary convergence region and all its Fourier coefficients continue meromorphically.

Insert \eqref{pc:Tpmeromorphic} in $\check T_p(u)*T_q(v)$. Constant-distribution terms vanish when $r\ge1$, since at least one factor has zero mean and a positive derivative. The two cross terms and the delta--delta term give
\begin{align*}
 \pccC_{pq}(u,v)={}&(-1)^p\frac{P_p(u)\pccD^r\pccZ_{1+v}-\pccE(u,v)\pccD^r\pccZ_{1+w}}u\\
 &+(-1)^q\frac{P_q(v)\check{\pccD^r\pccZ_{1+u}}
                   -\pccE(v,u)\check{\pccD^r\pccZ_{1+w}}}v
 +\frac{(-1)^pP_p(u)P_q(v)}{uv}\delta_0^{(r)}.
\end{align*}
Replace $\pccD^r\pccZ_{1+t}$ by $T_r(t)-P_r(t)\delta_0^{(r)}/t$ and use $\check\delta_0^{(r)}=(-1)^r\delta_0^{(r)}$. The nonlocal terms are exactly \eqref{pc:Jfpgen}. The remaining coefficient of $\delta_0^{(r)}$, after factoring $(-1)^p$, is
\[
 -\frac{P_p(u)P_r(v)}{uv}-\frac{P_q(v)P_r(u)}{uv}
 +\frac{P_p(u)P_q(v)}{uv}
 +\frac{P_r(w)}w\left(\frac{\pccE(u,v)}u+\frac{\pccE(v,u)}v\right).
\]
Equation \eqref{pc:EA} gives \eqref{pc:Deltagen}.

For $p=q=0$, keep the zero-mode terms explicitly. Since
$T_0(u)=\pccZ_{1+u}+(\delta_0-1)/u$, their convolution is
\begin{align*}
 \pccC_{00}(u,v)={}&\frac{\pccZ_{1+v}-\pccE(u,v)\pccZ_{1+w}}u
 +\frac{\check\pccZ_{1+u}-\pccE(v,u)\check\pccZ_{1+w}}v
 +\frac{\delta_0-1}{uv}.
\end{align*}
Replace $\pccZ_{1+t}=T_0(t)+(1-\delta_0)/t$ and apply \eqref{pc:EA}. The result is
\[
 \pccC_{00}(u,v)=\pccJ_{00}(u,v)+\frac{\pccA(u,v)-1}{uv}\delta_0,
\]
which is exactly \eqref{pc:Deltagen} with all $P$'s equal to one.

The numerator \eqref{pc:N} vanishes identically on $u=0$ and on $v=0$, because $P_j(0)=1$ and $\pccA$ equals one on both axes. Therefore it is divisible by $uv$ as a holomorphic germ. Taking coefficients proves \eqref{pc:mainformula}. The displayed calculation exhausts the difference, and produces only $\delta_0^{(r)}$; hence there are no omitted lower contacts.
\end{proof}

\subsection{Why the correction has no Stieltjes constants}
The logarithm of the gamma--trigonometric factor is
\begin{align}\label{pc:logA}
 \log\pccA(u,v)={}&\sum_{k\ge2}\frac{\zeta(k)}k
               \bigl(u^k+v^k-(u+v)^k\bigr)\notag\\
 &-\sum_{j\ge1}\frac{(1-2^{-2j})\zeta(2j)}j
               \bigl((u-v)^{2j}-(u+v)^{2j}\bigr).
\end{align}
Euler's constant cancels from the gamma ratio. It follows that
\[
 \Delta_{mn}^{p,q}\in\Q[\zeta(2),\zeta(3),\ldots],
\]
with the finite harmonic data encoded in rational coefficients of $P_p,P_q,P_r$. In particular there are no ordinary or generalized Stieltjes constants and no spectral derivatives of zeta in the contact coefficient. This is an assertion about the explicit expression, not an arithmetic independence theorem.

The reflection law is
\begin{equation}\label{pc:reflection}
 \Delta_{nm}^{q,p}=(-1)^r\Delta_{mn}^{p,q}.
\end{equation}
It follows both from \eqref{pc:N} and from reflecting the convolution and its coordinate extension. Reflection changes $\delta_0^{(r)}$ by precisely the necessary sign.

\subsection{All derivative orders at Stieltjes index zero}
\begin{corollary}[Polygamma contact formula]\label{pc:polygamma}
For $r=p+q$, including $r=0$,
\begin{equation}\label{pc:polyformula}
 \boxed{\ \Delta_{00}^{p,q}=(-1)^p
 \left[\frac{\pi^2}{3}-H_r^{(2)}+(H_r-H_p)(H_r-H_q)\right].\ }
\end{equation}
\end{corollary}
\begin{proof}
The coefficient of $uv$ in $\pccA$ is $2\zeta(2)$, and its linear part is zero. The coefficient of $uv$ in $P_r(u+v)$ is $H_r^2-H_r^{(2)}$. Extracting the remaining products in \eqref{pc:N} gives
\[
 2\zeta(2)+H_r^2-H_r^{(2)}-(H_p+H_q)H_r+H_pH_q,
\]
which is \eqref{pc:polyformula}.
\end{proof}
Some low-order values are
\begin{center}
\begin{tabular}{@{}ccc@{}}
\toprule
$(p,q)$ & $r$ & $\Delta_{00}^{p,q}$\\
\midrule
$(0,0)$&0&$\pi^2/3$\\
$(0,1)$&1&$\pi^2/3-1$\\
$(1,0)$&1&$1-\pi^2/3$\\
$(1,1)$&2&$1-\pi^2/3$\\
$(0,2)$&2&$\pi^2/3-5/4$\\
$(1,2)$&3&$13/12-\pi^2/3$\\
\bottomrule
\end{tabular}
\end{center}
These are corrections to a \emph{distribution in the shift}, not the values of the common-point product integrals studied in \cite{pcc:reports}.

\subsection{A closed parity law and derivative-split stabilization}
\begin{corollary}[The complete first row]\label{pc:edge}
For every $m\ge0$,
\begin{equation}\label{pc:edgeformula}
 \Delta_{m0}^{0,0}=m!\zeta(m+2)
 \begin{cases}
 3-2^{-m},&m\text{ even},\\
 1,&m\text{ odd}.
 \end{cases}
\end{equation}
\end{corollary}
\begin{proof}
Differentiating \eqref{pc:A} in $v$ and setting $v=0$ gives
\begin{equation}\label{pc:Aedge}
 \partial_v\pccA(u,0)=\psi(1-u)-\psi(1)+\pi\tan\frac{\pi u}{2}.
\end{equation}
The coefficient of $u^k$ is $-\zeta(k+1)$ for even $k\ge2$, and
$(3-2^{1-k})\zeta(k+1)$ for odd $k\ge1$. This follows from the gamma series and the cosine product, or by differentiating \eqref{pc:logA}. Substitute $k=m+1$ in \eqref{pc:deltacoeff}.
\end{proof}
For example,
\begin{align*}
 \Delta_{10}^{0,0}&=\zeta(3),&
 \Delta_{20}^{0,0}&=\tfrac{11}{2}\zeta(4),&
 \Delta_{30}^{0,0}&=6\zeta(5),\\
 \Delta_{11}^{0,0}&=\tfrac72\zeta(4),&
 \Delta_{21}^{0,0}&=4\zeta(5)+4\zeta(2)\zeta(3),&
 \Delta_{22}^{0,0}&=4\zeta(3)^2+\tfrac{83}{2}\zeta(6).
\end{align*}

\begin{corollary}[Stabilization of the derivative split]\label{pc:stabilization}
If $m\ge r$ or $n\ge r$, then
\begin{equation}\label{pc:stabilizationformula}
 \Delta_{mn}^{p,q}=(-1)^p\Delta_{mn}^{0,r}.
\end{equation}
\end{corollary}
\begin{proof}
Every product in \eqref{pc:N} other than $\pccA P_r(u+v)$ has degree at most $r$ in each variable. The indicated coefficient therefore vanishes in those products as soon as either $m+1>r$ or $n+1>r$. The remaining expression depends only on $r$, apart from the factor $(-1)^p$.
\end{proof}
An explicit first-row generator for arbitrary derivative orders is also useful:
\begin{align}\label{pc:general-edge}
 \Delta_{m0}^{p,q}=(-1)^{m+p}m![u^{m+1}]\bigl[
 &P_r(u)\partial_v\pccA(u,0)+P_r'(u)\notag\\
 &-H_qP_r(u)+(H_q-H_r)P_p(u)\bigr].
\end{align}
For $m\ge r$, the last three polynomial terms disappear. Formula \eqref{pc:Aedge} then gives a finite sum of zeta values with indices from $m+2-r$ through $m+2$, weighted by elementary symmetric polynomials in reciprocal integers.

\section{Exact convergent integral identities from finite Fourier data}\label{pc:weighted}
The distributional result produces ordinary integrals once the test function vanishes to sufficient order at the collision. Put
\begin{equation}\label{pc:weight}
 V_{M,k}(a)=(1-\cos(2\pi k a))^M,
 \qquad M,k\in\N_{>0}.
\end{equation}
Its finite Fourier expansion is
\begin{equation}\label{pc:weightfourier}
 V_{M,k}(a)=2^{-M}\sum_{j=-M}^M(-1)^j
              \binom{2M}{M+j}e^{2\pi\iu jk a}.
\end{equation}
This follows by expanding $2^{-M}(2-e^{\iu\theta}-e^{-\iu\theta})^M$. It vanishes to order $2M$ at both identified endpoints.

\begin{theorem}[Ordinary all-index correlation integrals]\label{pc:weightedtheorem}
If $2M>r=p+q$, then
\begin{align}\label{pc:weightedformula}
 \int_0^1 V_{M,k}(a)J_{mn}^{p,q}(a)\dd a
 =2^{1-M}\sum_{j=1}^M(-1)^j\binom{2M}{M+j}
 \Rea\!\left(\widehat T_{m,p}(-jk)\widehat T_{n,q}(jk)\right).
\end{align}
The integral is ordinary and absolutely convergent. The right-hand side is a finite exact expression given by \eqref{pc:Tfourier}.
\end{theorem}
\begin{proof}
The endpoint expansion in Section~\ref{pc:offpoint} bounds the weighted integrand by a constant times
$a^{2M-r-1}(1+|\log a|^{m+n+1})$ near zero, and similarly at one. Since $2M>r$, the integral converges. All derivatives of $V_{M,k}$ through degree $r$ vanish at zero, so the contact in \eqref{pc:mainformula} pairs to zero. Pair the convolution distribution with \eqref{pc:weightfourier}; its zero Fourier coefficient is zero. Realness of the original distributions combines the positive and negative modes into the displayed real part.
\end{proof}
At Stieltjes index zero, put $\lambda_j=\gamma+\log(2\pi jk)$. Then the Fourier product in \eqref{pc:weightedformula} is
\begin{equation}\label{pc:polyfourierproduct}
 (-1)^p(2\pi\iu jk)^r
 \left[(\lambda_j-H_p)(\lambda_j-H_q)+\frac{\pi^2}{4}
                  +\frac{\pi\iu}{2}(H_q-H_p)\right].
\end{equation}
For even $r=2d$ its real part is
\[
 (-1)^{p+d}(2\pi jk)^r
 \left[(\lambda_j-H_p)(\lambda_j-H_q)+\frac{\pi^2}{4}\right],
\]
whereas for odd $r=2d+1$ it is
\[
 (-1)^{p+d+1}(2\pi jk)^r\frac\pi2(H_q-H_p).
\]
Thus all logarithms disappear from these symmetric weighted polygamma correlations when the total derivative degree is odd.

At the first index, with $\lambda_k=\gamma+\log(2\pi k)$,
\begin{equation}\label{pc:Jweightedexample}
 \int_0^1(1-\cos(2\pi ka))
 \left[\gamma_1(a)+\gamma_1(1-a)-\frac{\pi^2}{3}\right]\dd a
 =-\lambda_k^2-\frac{\pi^2}{4}.
\end{equation}
The corresponding single-factor consequences of \eqref{pc:Tfourier} are
\begin{align}
 \int_0^1(1-\cos(2\pi kx))[-\psi(x)]\dd x&=\lambda_k,\label{pc:single0}\\
 \int_0^1(1-\cos(2\pi kx))\gamma_1(x)\dd x
 &=\frac{\pi^2}{24}-\frac12\lambda_k^2.\label{pc:single1}
\end{align}
These elementary Fourier specializations are included as checks and applications of the conventions, not as claims of priority over the classical Fourier theory.

\section{A Hurwitz kernel for all positive-integer spectral jets}\label{pc:Hurwitzweighted}
A second route makes the spectral zeta derivatives in the weighted identities explicit. Define the finite exponential sum
\begin{equation}\label{pc:SM}
 S_M(t)=\sum_{j=1}^M(-1)^j\binom{2M}{M+j}j^t,
 \qquad j^t=e^{t\log j}.
\end{equation}
It is entire, and $S_M^{(h)}(t)$ is the same finite sum with a factor $\log^h j$.

\begin{theorem}[Weighted Hurwitz kernel]\label{pc:Hurwitzkernel}
For $\Rea s<2M+1$, $s\ne1$,
\begin{equation}\label{pc:HurwitzkernelF}
 \int_0^1 V_{M,k}(x)\zeta(s,x)\dd x
 =2^{1-M}\Gamma(1-s)(2\pi k)^{s-1}
                  \sin\frac{\pi s}{2}\,S_M(s-1).
\end{equation}
At $s=2,\ldots,2M$, the right-hand side is interpreted by its removable limit. The identity can be spectrally differentiated to every finite order in this domain, except at the genuine pole $s=1$.
\end{theorem}
\begin{proof}
For $\Rea s<0$, pair the absolutely convergent Hurwitz Fourier expansion with \eqref{pc:weightfourier}. Its zero mode vanishes. The two signs of each nonzero frequency combine to
$2\Gamma(1-s)(2\pi jk)^{s-1}\cos(\pi(s-1)/2)$, and the cosine equals $\sin(\pi s/2)$. This gives \eqref{pc:HurwitzkernelF}.

For the larger domain, use $\zeta(s,x)=x^{-s}+\zeta(s,1+x)$. The weight is $O(x^{2M})$, so the first term and each of its spectral derivatives are integrable when $\Rea s<2M+1$. The shifted factor is holomorphic in $s$ except for its usual pole at one and locally bounded in $x$. Analytic continuation proves the formula on the stated connected domain and shows that the apparent poles at the other displayed integers are removable.
\end{proof}
The cancellations are finite identities, not numerical limiting phenomena:
\begin{equation}\label{pc:SMzeros}
 S_M(2d)=0\qquad(1\le d<M).
\end{equation}
Indeed, differentiate \eqref{pc:weightfourier} $2d$ times at zero and use the order-$2M$ zero of the weight. At even spectral integers the sine factor instead cancels the gamma pole.

\subsection{Complete holomorphic resonance germs}
The next formulas give every derivative, not just the value at a resonance. Let $h$ be a small spectral perturbation.
For an odd integer $s_0=2d+1$, with $1\le d<M$,
\begin{align}\label{pc:oddgerm}
 \int_0^1V_{M,k}(x)\zeta(2d+1+h,x)\dd x
 ={}&\frac{(-1)^{d+1}2^{1-M}(2\pi k)^{2d}}{(2d)!}\notag\\
 &\times\frac{\Gamma(1-h)(2\pi k)^h\cos(\pi h/2)}{P_{2d}(h)}
             \frac{S_M(2d+h)}h.
\end{align}
For an even integer $s_0=2d$, with $1\le d\le M$,
\begin{align}\label{pc:evengerm}
 \int_0^1V_{M,k}(x)\zeta(2d+h,x)\dd x
 ={}&\frac{(-1)^d2^{1-M}(2\pi k)^{2d-1}}{(2d-1)!}\notag\\
 &\times\frac{\Gamma(1-h)(2\pi k)^h}{P_{2d-1}(h)}
          \frac{\sin(\pi h/2)}h S_M(2d-1+h).
\end{align}
Both right-hand sides are holomorphic. These formulas follow directly from
\[
 \Gamma(-n-h)=\frac{(-1)^{n+1}\Gamma(1-h)}{n!\,hP_n(h)}.
\]
Consequently, the integral with $\zeta^{(j)}(s_0,x)$ is exactly $j!$ times the coefficient of $h^j$ in the corresponding right-hand side. The logarithm of its common gamma--harmonic factor is
\begin{equation}\label{pc:jetfactor}
 \log\frac{\Gamma(1-h)(2\pi k)^h}{P_n(h)}
 =(\gamma+\log(2\pi k)-H_n)h
 +\sum_{j\ge2}\frac{\zeta(j)+(-1)^jH_n^{(j)}}j h^j.
\end{equation}
This is a finite coefficient algorithm in gamma, harmonic numbers, ordinary zeta values, and logarithms of the integers $1,\ldots,M$. No derivative of an infinite special-function series has to be fitted numerically.

At the pole $s=1$, put $c_M=2^{-M}\binom{2M}{M}$, the mean of the weight. The Stieltjes counterpart is the holomorphic identity
\begin{align}\label{pc:stieltjesweightedgerm}
 \int_0^1V_{M,k}(x)g(h,x)\dd x
 =\frac{-2^{1-M}\Gamma(1-h)(2\pi k)^h\cos(\pi h/2)S_M(h)-c_M}{h}.
\end{align}
Its numerator vanishes at zero because $S_M(0)=-\binom{2M}{M}/2$. Extracting $(-1)^m m![h^m]$ evaluates the ordinary weighted integral of $\gamma_m(x)$ at every index. For $M=1$, this recovers \eqref{pc:single0} and \eqref{pc:single1}.

\subsection{Explicit logarithmic identities at order three}
For every positive integer $k$, put $\lambda_k=\gamma+\log(2\pi k)$. Taking $M=2,d=1$ in \eqref{pc:oddgerm}, where $S_2(t)=-4+2^t$, gives
\begin{align}
 \boxed{\ \int_0^1(1-\cos(2\pi kx))^2\zeta(3,x)\dd x
       =4\pi^2k^2\log2.\ }\label{pc:zeta3}\\[4pt]
 \boxed{\ \int_0^1(1-\cos(2\pi kx))^2\zeta'(3,x)\dd x
       =4\pi^2k^2\log2
          \left(\lambda_k-\frac32+\frac{\log2}{2}\right).\ }\label{pc:zeta3prime}
\end{align}
Here $\zeta'$ denotes the spectral derivative. To see the derivative constant explicitly, the common factor in \eqref{pc:oddgerm} has logarithmic derivative $\lambda_k-H_2=\lambda_k-3/2$ at zero, while
\[
 \frac{S_2(2+h)}h=4\log2+2\log^22\,h+O(h^2).
\]
The local integrands behave like $x$ and $x\log x$, respectively, so these are ordinary absolutely convergent integrals. Their polygamma form follows from $\psi^{(2)}(x)=-2\zeta(3,x)$.

The correlation identity in Theorem~\ref{pc:weightedtheorem}, at $(p,q)=(1,1)$, gives the compatible combination
\begin{align}\label{pc:correlation3}
 \int_0^1(1-\cos(2\pi kx))^2
       \bigl[\zeta(3,x)+2\zeta'(3,x)\bigr]\dd x
 =4\pi^2k^2\bigl[2\log2(\lambda_k-1)+\log^22\bigr].
\end{align}
This agreement checks the derivative signs and the harmonic constants by two exact routes: the shifted-contact calculus and the single-Hurwitz spectral kernel.

\section{Polylogarithm boundary values and invisible local terms}\label{pc:polylogs}
For $\sigma\in\{+1,-1\}$, let
\begin{equation}\label{pc:periodicLi}
 \mathfrak L_s^\sigma(a)=\sum_{k\ge1}\frac{e^{\sigma2\pi\iu ka}}{k^s}
\end{equation}
be interpreted as a periodic distribution. Its Fourier coefficients have polynomial growth on every compact spectral set, so this is an entire distribution-valued function of $s$. It is also the Abel limit of $\Li_s(\rho e^{\sigma2\pi\iu a})$ as $\rho\uparrow1$. Spectral derivatives multiply the coefficients by powers of $-\log k$. Off the integer points these distributions agree with the usual analytic boundary values of polylogarithms \cite{rw:DLMFPolylog}.

Put $c=\gamma+\log(2\pi)$ and expand the finite polynomial
\begin{equation}\label{pc:polylogpolynomial}
 (-1)^p(\sigma2\pi\iu)^r
 \bigl[b_m(c+X-\sigma\iu\pi/2)+\kappa_{m,p}\bigr]
 \bigl[b_n(c+X+\sigma\iu\pi/2)+\kappa_{n,q}\bigr]
 =\sum_{j=0}^{m+n+2}a_{j,\sigma}X^j.
\end{equation}
Equation \eqref{pc:Tfourier} proves the all-index polylogarithmic identity
\begin{equation}\label{pc:polylogidentity}
 \pccC_{mn}^{p,q}=
 \sum_{\sigma=\pm1}\sum_{j=0}^{m+n+2}
 a_{j,\sigma}(-1)^j
 \left.\partial_s^j\mathfrak L_s^\sigma\right|_{s=-r}.
\end{equation}
This is a finite formula; its equality is checked at every Fourier mode, including the zero mode, which vanishes on both sides. Together with Theorem~\ref{pc:main}, it identifies the exact delta derivative that must be added to the coordinate extension of the corresponding off-point polylogarithm identity.

The elementary cautionary example is
\begin{equation}\label{pc:Li0delta}
 \mathfrak L_0^++\mathfrak L_0^-=\delta_0-1.
\end{equation}
For $a\notin\Z$, the ordinary rational values satisfy
$\Li_0(e^{2\pi\iu a})+\Li_0(e^{-2\pi\iu a})=-1$. The missing delta cannot be read from that pointwise identity. Similarly, for $r\ge1$,
\begin{equation}\label{pc:Lirdelta}
 \mathfrak L_{-r}^++(-1)^r\mathfrak L_{-r}^-
 =\frac{\delta_0^{(r)}}{(2\pi\iu)^r}.
\end{equation}
These identities follow immediately by adding positive and negative Fourier modes, or by differentiating \eqref{pc:Li0delta}. They explain why substituting rational negative-order polylogarithms before fixing distributional boundary values can erase precisely the terms computed by the contact theorem.

\section{Exact primitives and the log-gamma bridge}\label{pc:primitives}
Define the mean-zero periodic inverse derivative by
\[
 \widehat{\pccI U}(k)=\frac{\widehat U(k)}{2\pi\iu k}\quad(k\ne0),
 \qquad\widehat{\pccI U}(0)=0.
\]
It satisfies $\pccD\pccI U=U-\widehat U(0)1$. Let $B_h$ be the Bernoulli polynomial, periodically interpreted when necessary. Its standard Fourier coefficients, obtainable also from the negative-integer Hurwitz values \cite{intdist:DLMFHurwitz}, give
\begin{equation}\label{pc:primitivecontact}
 \pccI^h\delta_0^{(r)}=
 \begin{cases}
 \delta_0^{(r-h)},&h<r,\\
 \delta_0-1,&h=r\ge1,\\
 -B_{h-r}(\{a\})/(h-r)!,&h>r.
 \end{cases}
\end{equation}
For $h=r=0$, the left side is simply $\delta_0$ and the middle rule is not used.

\begin{corollary}[All normalized contact primitives]\label{pc:primitivecor}
Apply $\pccI^h$ to Theorem~\ref{pc:main}. For $h>r$ the contact correction is exactly
\begin{equation}\label{pc:bernoulli}
 \pccI^h\pccC_{mn}^{p,q}-\pccI^h\pccJ_{mn}^{p,q}
 =-\Delta_{mn}^{p,q}\frac{B_{h-r}(\{a\})}{(h-r)!}.
\end{equation}
There is no undetermined polynomial or integration constant: the zero Fourier mode is fixed by definition.
\end{corollary}
For $h\ge r+2$, the integrated convolution has an absolutely convergent Fourier series because its coefficients are
$O(|k|^{r-h}(1+\log^{m+n+2}|k|))$. Thus sufficiently integrated identities are equalities of continuous ordinary functions, even though their derivatives contain contact distributions.

For a concrete gamma example, set
\[
 \ell(x)=\log\Gamma(x)-\frac12\log(2\pi).
\]
Raabe's mean and the Hurwitz identity for $\zeta'(0,x)$ give $\int_0^1\ell(x)\dd x=0$. Distributionally $\pccD\ell=-T_{0,0}$; no extra delta occurs because the constant term in the endpoint expansion of $\log\Gamma(x)$ is zero and $\log\Gamma(1)=0$. Hence $\ell=-\pccI T_{0,0}$ and
\begin{equation}\label{pc:loggammacov}
 K_\Gamma(a):=\int_0^1\ell(x)\ell(\{x+a\})\dd x
 =-\pccI^2\pccC_{00}^{0,0}.
\end{equation}
The integral exists also at $a=0$ because logarithmic endpoint singularities are square integrable. The contact theorem gives the normalized form
\begin{equation}\label{pc:gammaBernoulli}
 K_\Gamma(a)=-\pccI^2\pccJ_{00}^{0,0}(a)+\frac{\pi^2}{6}B_2(\{a\}).
\end{equation}
The Bernoulli term is the integrated image of the $\pi^2\delta_0/3$ term in \eqref{pc:preview}.

The ordinary Fourier representation is already
\eqref{shj:gammapoly}, after subtracting the Raabe mean. The additional
information here is its fixed-coordinate Bernoulli correction: it is the
integrated image of the collision contact, which a punctured-circle
polylogarithm identity cannot determine.

\paragraph{Evidence and remaining scope.}
The two independent contact generators agree after aligning their
derivative signs and unit-coordinate conventions. Their finite exact
tables and independent subtracted Fourier-mode quadratures test those
normalizations; the analytic comparison above proves the all-index law.
This completes the specified periodic extension problem. Cubic geometric
collision limits, ordered higher-depth germs and canonical changes of
subtraction coordinates retain their separate audit status. The contact
coefficient does not assert independence of its displayed zeta coordinates.

